% Generated from article.tex by build_standalone.py.
% Edit the modular source and regenerate this file.

% BEGIN preamble.tex
\documentclass[11pt]{article}
\usepackage[T1]{fontenc}
\usepackage{lmodern}
\usepackage[margin=1in]{geometry}
\usepackage{amsmath,amssymb,amsthm,mathtools,mathrsfs}
\usepackage{booktabs,longtable,tabularx,array,enumitem,xcolor,microtype,xurl}
\usepackage{graphicx,needspace}
\definecolor{inkblue}{RGB}{24,61,97}
\usepackage[colorlinks=true,linkcolor=inkblue,citecolor=inkblue,urlcolor=inkblue]{hyperref}
\usepackage{bookmark}
\hypersetup{pdftitle={Barnes Resolvents and Higher Harmonic Identities},
 pdfauthor={Research continuation prepared for Vladimir Reshetnikov},
 pdfsubject={Exact polylogarithm, Gamma, harmonic, Stieltjes, and Tornheim identities},
 pdfkeywords={polylogarithm, Barnes G, multiple gamma, polygamma, Hurwitz zeta, Stieltjes constants, Tornheim}}
\setcounter{tocdepth}{2}
\numberwithin{equation}{section}
\allowdisplaybreaks[2]
\emergencystretch=3em
\setlength{\parskip}{3pt}
\theoremstyle{plain}
\newtheorem{theorem}{Theorem}[section]
\newtheorem{proposition}[theorem]{Proposition}
\newtheorem{lemma}[theorem]{Lemma}
\newtheorem{corollary}[theorem]{Corollary}
\theoremstyle{definition}
\newtheorem{definition}[theorem]{Definition}
\newtheorem{example}[theorem]{Example}
\theoremstyle{remark}
\newtheorem{remark}[theorem]{Remark}
\DeclareMathOperator{\Li}{Li}
\DeclareMathOperator{\Log}{Log}
\DeclareMathOperator{\FP}{FP}
\DeclareMathOperator{\CT}{CT}
\DeclareMathOperator{\Res}{Res}
\DeclareMathOperator{\sgn}{sgn}
\newcommand{\C}{\mathbb C}
\newcommand{\R}{\mathbb R}
\newcommand{\N}{\mathbb N}
\newcommand{\Q}{\mathbb Q}
\newcommand{\Z}{\mathbb Z}
\newcommand{\ii}{\mathrm i}
\newcommand{\dd}{\,\mathrm dx}
\newcommand{\src}[1]{\nolinkurl{#1}}
\newcommand{\Liord}[2]{\Li_{#1}^{[1]}(#2)}
\newcommand{\snapshot}{\nolinkurl{7bd45777a8a127e5dc69aff8887ca36a79a3059d}}
\newcolumntype{Y}{>{\raggedright\arraybackslash}X}
\title{\textcolor{inkblue}{Barnes Resolvents and Higher Harmonic Identities}\\[8pt]
\large Polynomial Hurwitz transforms, complex kernel powers,\\
quartic Stieltjes moments, and transverse Tornheim data}
\author{Research continuation prepared for Vladimir Reshetnikov}
\date{11 October 2026}
\hypersetup{pdftitle={Barnes Resolvents and Higher Harmonic Identities},pdfauthor={Research continuation for Vladimir Reshetnikov}}

% END preamble.tex

\begin{document}
\maketitle
\begin{abstract}
We continue the exact-identity programme of the ProveIt polylogarithm
manuscript from a fixed repository snapshot. A polynomially weighted
Hurwitz transform gives the individual cotangent resolvents of every
unit-period multiple Gamma function with an explicit normalization.
For the Barnes $G$ function, the answer uses ordinary polylogarithms,
their order derivatives, and one convergent depth-two derivative; an
antisymmetric component reduces entirely to depth one. The same
transform gives all weighted Stieltjes coefficients and their full
endpoint corrections. We also develop complex powers of the kernel,
an exact quartic digamma--harmonic bridge, and an independently
convergent transverse coordinate for higher cyclic Tornheim
derivatives. Every reduction is stated relative to an explicit
function class; no arithmetic independence or literature-wide
priority claim is made. Proofs, normalization checks, reproducible
computations, a narrowly scoped correction to a published Barnes
identity, and further research questions accompany the results.
\end{abstract}
\tableofcontents
\newpage

% BEGIN sections/01_scope.tex
\section{Scope, source baseline, and conventions}
\label{bhr:sec:scope}

The starting point is the ProveIt repository at commit
\snapshot, whose recorded time is 11 October 2026, 01:57:13 UTC.
We obtained the canonical manuscript's 81 chapter files, its README,
validation record, preamble and references, and the text members of all
eleven archives then present in \src{docs/incoming}. The source review
was targeted at the identities used below; it was not a line-by-line
soundness audit of every volume. The accompanying source inventory
identifies the snapshot and the incoming Git blob identifiers.

The newest source, \emph{Harmonic Parity and Resolvent Identities}
\cite{bhr:Parity}, already settles the rational cotangent-resolvent
problem, the balanced negapolygamma ladder, and a quartic directional
normal form. Those conclusions are baseline results here. Its
questions R4 and R6 instead ask for nonintegral complex kernel powers
and individual multiple Gamma transforms. The cubic harmonic report
\cite{bhr:Cubic} asks for a further bridge from higher digamma powers to
harmonic Laurent coefficients. These are the principal targets of
the present continuation.

\begin{table}[htbp]
\small
\begin{tabularx}{\textwidth}{@{}>{\raggedright\arraybackslash}p{0.23\textwidth}Y@{}}
\toprule
Source target & Exact scope pursued in this article\\
\midrule
Parity R6: individual Barnes functions & A finite formula for every
polynomially weighted, unit-period multiple Gamma resolvent, including
all normalization polynomials; an explicit $G$ formula.\\[5pt]
Parity R4: complex powers & A joint spectral transform and conversion
at positive integer exponent crossings, with the branch fixed along
the integration interval.\\[5pt]
Cubic Q6: higher digamma powers & The fourth-power finite part in
ordinary zeta/Stieltjes data and explicitly defined depth-two and
depth-three harmonic Laurent coefficients; an absolutely convergent
series representation.\\[5pt]
Parity R10: a transverse cyclic coordinate & A convergent integral
representative and its use in the fifth cyclic directional layer.\\
\bottomrule
\end{tabularx}
\caption{Developments relative to the inspected reports. Theorems below
state the precise domains; a reduction to higher-depth data is not a
reduction to an unspecified algebra of familiar constants.}
\label{bhr:tab:targets}
\end{table}

\subsection{What remains a conjecture}

The short Gaussian $S_6$ evaluation with label
\src{cycloquot:conj:S6} and the revised $S_8$ evaluation with label
\src{s8new:conj:S8} remain conjectural. The older frozen $S_8$ vector
was rigorously rejected in the existing manuscript; it is a different
candidate. None of the transformations in this article is an exact
certificate for either current candidate. Similarly, the familiar
cubic constants are not declared to belong to a smaller classical
constant algebra simply because an integral or a new convergent
series represents them.

\subsection{Notation}

We use the standard Hurwitz Laurent convention
\begin{equation}\label{bhr:eq:stieltjes-convention}
 \zeta(1-s,x)=-\frac1s+
    \sum_{m=0}^{\infty}\gamma_m(x)\frac{s^m}{m!},
 \qquad \gamma_0(x)=-\psi(x),\quad \psi=\frac{\Gamma'}\Gamma.
\end{equation}
Thus $\gamma_m=\gamma_m(1)$ and $\gamma=\gamma_0$.
An ordinary prime on $\zeta$ means a derivative in its first,
spectral argument. A superscript on $\psi$ or $\gamma_m(x)$ means
an argument derivative. We write
\[
 \Li_s^{[j]}(q)=\partial_s^j\Li_s(q),\qquad
 H_n^{(r)}=\sum_{k=1}^n k^{-r},\qquad H_n=H_n^{(1)}.
\]
All logarithms of positive real numbers are real. Complex logarithms
are fixed explicitly whenever they occur in a power. Products of
positive real Gamma and Barnes functions on $(0,1)$ use their real
logarithms.

The endpoint finite part in this article is
\begin{equation}\label{bhr:eq:fp-definition}
 \FP_x\int_0^1 f(x)\,dx
   = [\delta^0(\log\delta)^0]
       \int_{\delta}^{1-\delta}f(x)\,dx.
\end{equation}
This definition applies when the integral has a finite asymptotic
expansion of nonvanishing powers and logarithms plus a constant and a
vanishing remainder. For complex exponents, the removed powers may
oscillate as well as diverge. Our proofs either establish that expansion or
give an equivalent explicit subtraction. A parameter Laurent
constant, denoted $\CT$, is a different operation. Their equality
is asserted only after the relevant correction has been proved.

\subsection{Classical inputs and the meaning of a contribution}

The Hurwitz Fourier formula, the argument-shift law, Gamma
derivatives, Bernoulli polynomial evaluations, and ordinary
polylogarithm calculus are standard \cite{bhr:DLMFZ,bhr:DLMFG}.
Espinosa--Moll already developed Hurwitz integrals and balanced
negapolygamma functions \cite{bhr:EMIntegrals,bhr:EMGeneralized}.
The multiple-Gamma/Hurwitz-zeta conversion is classical
\cite{bhr:Adamchik,bhr:Vardi}. The Fourier coefficients of
$x\log\Gamma(x)$ and $\log G(x)$ are also known, in particular in
Mez\H{o}'s paper \cite{bhr:Mezo}. We do not reclassify any of those
facts as discoveries.

Our Barnes result is a finite resolvent reduction, valid with an
arbitrary polynomial weight and every specified normalization
constant. Its spectral form supplies an all-index Stieltjes
extension at the same time. The later sections isolate the precise
additional harmonic and transverse coordinates rather than hiding
them inside unnamed finite parts. These are developments relative
to the pinned source programme. Numerical checks in the package
support independent error detection; the proofs in the text carry
the mathematical claims.

% END sections/01_scope.tex


% BEGIN sections/02_weighted_transform.tex
\section{A polynomially weighted Hurwitz resolvent}
\label{bhr:sec:weighted}

\subsection{The two Fourier directions and a convergent depth-two function}

For $z\in\C\setminus\R$, put
\begin{equation}\label{bhr:eq:coordinates}
 \begin{gathered}
 \epsilon=\sgn(\Im z),\quad \omega=2\epsilon\pi\ii,
 \quad \sigma=-\omega,\quad a=z+\epsilon\pi\ii,\\
 q=\frac{z-\epsilon\pi\ii}{z+\epsilon\pi\ii},\qquad
 c=\frac{2\epsilon\pi\ii}{a^2},\qquad |q|<1,\\
 \Log\omega=\log(2\pi)+\epsilon\pi\ii/2,\qquad
 \Log\sigma=\log(2\pi)-\epsilon\pi\ii/2.
 \end{gathered}
\end{equation}
Let $K(x)=\pi\cot(\pi x)$ and $R_z(x)=(K(x)-z)^{-1}$.
The elementary identity
\begin{equation}\label{bhr:eq:resolvent-fourier}
 R_z(x)=-\frac1a+c\sum_{n\ge1}q^{n-1}e^{\omega nx}
\end{equation}
converges uniformly with every argument derivative. In particular
$R_z(0)=R_z(1)=0$ and $R_z(x)=x+O(x^2)$ at zero.

The natural depth-two function for multiplication by a polynomial is
\begin{equation}\label{bhr:eq:omega-function}
 \mathcal O_{\alpha,\beta}(q)
   =\sum_{n,m\ge1}\frac{q^n}{(m+n)^\alpha m^\beta}
   =\Li_{\alpha,\beta}(q,q^{-1}).
\end{equation}
Here our convention is
$\Li_{\alpha,\beta}(u,v)=\sum_{k>m\ge1}
u^kv^m/(k^\alpha m^\beta)$.
The last expression in \eqref{bhr:eq:omega-function} is interpreted by
the displayed convergent series, not by assigning separate
independent branches to the two arguments.

\begin{lemma}[Normal convergence]\label{bhr:lem:omega-normal}
The series \eqref{bhr:eq:omega-function}, with all finite order
derivatives, converges normally on
\[
 |q|<1,\qquad \Re(\alpha+\beta)>1.
\]
It vanishes at $q=0$, and division by $q$ has a removable value there.
\end{lemma}
\begin{proof}
On a compact parameter set split the inner sum at $m=2n$.
For $m>2n$, the summand is bounded by a fixed multiple of
$|q|^n m^{-\Re(\alpha+\beta)}$. For $m\le2n$, all possible
growth is bounded by a fixed power of $n$, and the geometric
factor dominates it. Derivatives introduce only powers of
$\log m$, $\log(m+n)$ and $n$, which are handled by a small
margin in the strict inequalities. The same bounds justify
termwise differentiation. Every term contains $q^n$ with $n\ge1$.
\end{proof}

We shall also need the normally convergent subtracted functions
\begin{align}
 U_j(s,q)&=\sum_{n,m\ge1}q^n m^{-s}
       \bigl((m+n)^{-j}-m^{-j}\bigr),\label{bhr:eq:U}\\
 V_j(s,q)&=\sum_{n,m\ge1}q^n m^{-j}
       \bigl((m+n)^{-s}-m^{-s}\bigr),\label{bhr:eq:V}
\end{align}
where $j\ge1$. Each is holomorphic for $\Re s>-j$, $|q|<1$.
In the common initial domain they satisfy
\begin{equation}\label{bhr:eq:UV-continuation}
 U_j=\mathcal O_{j,s}-\frac{q}{1-q}\zeta(s+j),\qquad
 V_j=\mathcal O_{s,j}-\frac{q}{1-q}\zeta(s+j).
\end{equation}
For example, in the tail $m>2n$ each difference gains a factor
$O(n/m)$, uniformly on compact sets. The remaining part is
geometrically summable as in Lemma~\ref{bhr:lem:omega-normal}.
This proves the claimed larger domain directly, without assigning
values to divergent unsubtracted summands.

\subsection{The finite polynomial data}

For a polynomial $P$ of degree $d$, define
\begin{equation}\label{bhr:eq:polynomial-data}
 \mu_P=\int_0^1P(x)\,dx,\qquad
 b_j(P)=\frac{(-1)^{j-1}}{\omega^j}
    \bigl(P^{(j-1)}(1)-P^{(j-1)}(0)\bigr),\quad 1\le j\le d.
\end{equation}
The sum is empty when $P$ is constant. Repeated integration by
parts gives
\begin{equation}\label{bhr:eq:P-fourier}
 \int_0^1P(x)e^{\omega kx}\,dx
 =\begin{cases}
   \mu_P,& k=0,\\[2pt]
   \displaystyle\sum_{j=1}^d\frac{b_j(P)}{k^j},& k\in\Z\setminus\{0\}.
  \end{cases}
\end{equation}
The derivative of order $d$ is constant and hence contributes no
endpoint difference. Thus only these finitely many data occur.

Define the elementary polynomial resolvent
\begin{equation}\label{bhr:eq:Tpoly}
 T_P(z)=\int_0^1P(x)R_z(x)\,dx
       =-\frac{\mu_P}{a}+
          \frac cq\sum_{j=1}^d b_j(P)\Li_j(q).
\end{equation}
All quotients by $q$ in this article are given their removable
values at $q=0$, equivalently $z=\epsilon\pi\ii$.

\subsection{The master identity}

Set
\begin{align}\label{bhr:eq:Bpoly}
 \mathcal B_P(s,q)
   ={}&\sigma^{-s}\mu_P\Li_s(q)\notag\\
    &+\sum_{j=1}^d b_j(P)
       \left\{\omega^{-s}U_j(s,q)
        +\sigma^{-s}\bigl(\Li_s(q)\Li_j(q)
                              +(-1)^jV_j(s,q)\bigr)\right\}.
\end{align}

\begin{theorem}[Polynomial Hurwitz resolvent]\label{bhr:thm:weighted-master}
For every polynomial $P$, every $z\notin\R$, and $\Re s>0$,
\begin{equation}\label{bhr:eq:weighted-master}
 \boxed{\displaystyle
 I_P(s,z):=\int_0^1P(x)\zeta(1-s,x)R_z(x)\,dx
          =\frac{c\Gamma(s)}q\mathcal B_P(s,q).}
\end{equation}
The bracket $\mathcal B_P$ is holomorphic on $\Re s>-1$.
The formula therefore gives a meromorphic continuation through
$s=0$, with residue $-T_P(z)$ and no other pole in that half-plane.
On $\Re s>0$ it is a finite depth-two polylogarithmic formula
by \eqref{bhr:eq:UV-continuation}.
\end{theorem}

\begin{proof}
Initially take $\Re s>1$. The classical Hurwitz Fourier expansion
\cite[\S25.11]{bhr:DLMFZ} is absolutely convergent and reads
\begin{equation}\label{bhr:eq:hurwitz-fourier}
 \zeta(1-s,x)=\Gamma(s)\sum_{m\ge1}m^{-s}
       \left(\omega^{-s}e^{\omega mx}
                         +\sigma^{-s}e^{-\omega mx}\right).
\end{equation}
Multiply this by \eqref{bhr:eq:resolvent-fourier} and use
\eqref{bhr:eq:P-fourier}. The constant Fourier term of $R_z$ gives
\[
 -\frac{\Gamma(s)}a\sum_j b_j(P)
        \bigl(\omega^{-s}+(-1)^j\sigma^{-s}\bigr)\zeta(s+j).
\]
The equal-frequency diagonal contributes
$c\Gamma(s)\sigma^{-s}\mu_P\Li_s(q)/q$.
The same-sign frequencies contribute $\mathcal O_{j,s}$.
For the remaining frequencies split at $m=n$:
\begin{align}\label{bhr:eq:lattice-splitting}
 \sum_{\substack{n,m\ge1\\m\ne n}}
    \frac{q^n}{m^s(n-m)^j}
 &=\Li_s(q)\Li_j(q)+(-1)^j\mathcal O_{s,j}(q).
\end{align}
Indeed $m<n$ is parameterized by $n=m+h$ and gives the
product; $m>n$ is parameterized by $m=n+h$ and gives the second
sum. Absolute convergence justifies these changes of variables.

Insert \eqref{bhr:eq:UV-continuation}. Every ordinary zeta term
cancels because
\begin{equation}\label{bhr:eq:coefficient-cancellation}
 \frac{c}{1-q}=\frac1a.
\end{equation}
What remains is exactly \eqref{bhr:eq:weighted-master}.
Both sides are holomorphic on $\Re s>0$, so the initial identity
extends there. On the left, use
$\zeta(1-s,x)=x^{s-1}+\zeta(1-s,x+1)$:
the first term times $P R_z$ is locally integrable for $\Re s>-1$,
and the second has just the constant Hurwitz pole $-1/s$.
The normal convergence of $U_j,V_j$ proves the same assertion
on the right. Finally, telescoping in $m$ gives
\[
 U_j(0,q)=-\frac{\Li_j(q)}{1-q},\qquad V_j(0,q)=0.
\]
Since $\Li_0(q)=q/(1-q)$,
\begin{equation}\label{bhr:eq:Bzero}
 \mathcal B_P(0,q)=\mu_P\frac{q}{1-q}
                     -\sum_jb_j(P)\Li_j(q)=-\frac q c T_P(z).
\end{equation}
This proves the residue, including the removable $q=0$ case.
\end{proof}

\begin{example}[The first nonconstant weight]\label{bhr:ex:xweight}
For $P(x)=x$, one has $\mu_P=1/2$, $b_1=1/\omega$ and
\begin{align}
 \mathcal B_x(s,q)
  =\frac12\sigma^{-s}\Li_s(q)
   +\frac1\omega\left\{\omega^{-s}U_1(s,q)
       +\sigma^{-s}\bigl(\Li_s(q)\Li_1(q)-V_1(s,q)\bigr)\right\}.
 \label{bhr:eq:xweight}
\end{align}
This is the missing $x\log\Gamma(x)$ input to an individual
Barnes-$G$ transform. Setting $P=1$ deletes $U_j,V_j$ and
recovers the previously established unweighted formula.
\end{example}

\subsection{All weighted Stieltjes coefficients}

\begin{theorem}[A holomorphic Stieltjes generator]\label{bhr:thm:weighted-stieltjes}
For every polynomial $P$, integer $m\ge0$, and $z\notin\R$,
\begin{equation}\label{bhr:eq:weighted-stieltjes}
 \boxed{\displaystyle
 \int_0^1P(x)\gamma_m(x)R_z(x)\,dx
   =\frac{m!c}{q}[s^{m+1}]
               \Gamma(1+s)\mathcal B_P(s,q).}
\end{equation}
These are ordinary absolutely convergent integrals.
Equivalently their exponential generating function is
\begin{equation}\label{bhr:eq:weighted-stieltjes-germ}
 I_P(s,z)+\frac{T_P(z)}s
   =\frac cq\frac{\Gamma(1+s)\mathcal B_P(s,q)
                          -\mathcal B_P(0,q)}s.
\end{equation}
\end{theorem}
\begin{proof}
The shift law
$\gamma_m(x)=x^{-1}\log^m x+\gamma_m(x+1)$ and $R_z=O(x)$
prove convergence. The same bounds, locally uniformly for $s$
near zero, allow \eqref{bhr:eq:stieltjes-convention} to be integrated
after removing its constant pole. Now use Theorem~\ref{bhr:thm:weighted-master}
and \eqref{bhr:eq:Bzero}. The numerator on the right of
\eqref{bhr:eq:weighted-stieltjes-germ} vanishes at zero.
\end{proof}

For instance, the first two moments are obtained from
\begin{align}
 \int_0^1P\gamma_0R_z\,dx
  &=\frac cq\bigl(\partial_s\mathcal B_P(0,q)
                                      -\gamma\mathcal B_P(0,q)\bigr),\\
 \int_0^1P\gamma_1R_z\,dx
  &=\frac c{2q}\bigl(\partial_s^2\mathcal B_P(0,q)
       -2\gamma\partial_s\mathcal B_P(0,q)
                +(\gamma^2+\zeta(2))\mathcal B_P(0,q)\bigr).
\end{align}
Every derivative of the subtracted double sums is normally
convergent. In particular these formulas do not require subtracting
two divergent numerical evaluations of \eqref{bhr:eq:UV-continuation}.

\subsection{Argument derivatives and complete local corrections}

Let
\begin{equation}\label{bhr:eq:Ap}
 A_p(s)=\prod_{j=1}^p(s-j),\qquad
 r_{P,p}(z)=[x^p]\bigl(P(x)R_z(x)\bigr).
\end{equation}
The empty product is one. Meromorphic continuation of $I_P$ to
any desired finite spectral neighborhood can be constructed by
subtracting finitely many terms from the Taylor expansion of
$P(x)R_z(x)$ at zero. This supplies an unambiguous meaning to
the next formula even when $s-p$ is outside the normal domain
of the first subtraction in \eqref{bhr:eq:U}--\eqref{bhr:eq:V}.

\begin{theorem}[Weighted argument-derivative conversion]\label{bhr:thm:weighted-contact}
For $m\ge0$, $p\ge1$ and a polynomial $P$,
\begin{equation}\label{bhr:eq:weighted-contact}
 \boxed{\displaystyle
 \FP_x\int_0^1P(x)\gamma_m^{(p)}(x)R_z(x)\,dx
   =m![s^m]\left\{
       A_p(s)I_P(s-p,z)-\frac{r_{P,p}(z)A_p(s)}s\right\}.}
\end{equation}
The expression in braces is holomorphic near $s=0$.
The entire polynomial amplitude $A_p(s)$ must be subtracted,
not only its value at zero.
\end{theorem}
\begin{proof}
The Hurwitz argument ladder gives the meromorphic identity
\[
 \partial_x^p\zeta(1-s,x)=A_p(s)\zeta(1-s+p,x).
\]
The singular part at zero is $A_p(s)x^{s-p-1}$.
After multiplication by $P R_z$, the only exponent crossing
$-1$ at $s=0$ is $r_{P,p}A_p(s)x^{s-1}$.
Its continued integral is $r_{P,p}A_p(s)/s$, whereas the
finite part of every coefficient $x^{-1}\log^j x$ on $(0,1)$
is zero. All other terms commute with coefficient extraction.
To justify the latter statement, subtract a Taylor polynomial
through degree $p+1$ on $(0,\delta_0)$, integrate its terms
explicitly, and use normal integrability of the remainder.
The nonsingular Hurwitz tail at $x+1$ is analytic after an
argument derivative. This proves the formula and the
holomorphy assertion.
\end{proof}

For a completely explicit continuation of the depth-two expression,
one can subtract $N$ terms of $(1+n/m)^{-j}$ and $(1+n/m)^{-s}$.
The resulting known terms are respectively
\begin{align}
 &\sum_{k=0}^{N-1}\frac{(-1)^k(j)_k}{k!}
         \Li_{-k}(q)\zeta(s+j+k),\\
 &\sum_{k=0}^{N-1}\frac{(-1)^k(s)_k}{k!}
         \Li_{-k}(q)\zeta(s+j+k).
\end{align}
The remainders are normally convergent for $\Re s>1-j-N$.
This follows by the same $m>2n$ split and the finite Taylor
remainder estimate; the terms $m\le2n$ are still geometrically
summable. Thus \eqref{bhr:eq:weighted-contact} is a constructive
finite calculation plus convergent sums at every fixed $m,p$.

Repeated resolvents follow by $z$ differentiation:
\begin{equation}\label{bhr:eq:weighted-repeated}
 \FP_x\int_0^1\frac{P(x)\gamma_m^{(p)}(x)}{(K(x)-z)^r}\,dx
  =\frac1{(r-1)!}\partial_z^{r-1}
        \FP_x\int_0^1P(x)\gamma_m^{(p)}(x)R_z(x)\,dx.
\end{equation}
Local subtractions are holomorphic in $z$ on compact subsets of
either half-plane. If $P$ vanishes to order $e$ at zero, the
integral in \eqref{bhr:eq:weighted-repeated} is absolutely convergent
when $e+r>p$.

% END sections/02_weighted_transform.tex


% BEGIN sections/03_barnes.tex
\section{Individual Barnes functions and all multiple Gamma ranks}
\label{bhr:sec:barnes}

The balanced functions in the preceding reports are particular
linear combinations of Hurwitz-zeta derivatives and Bernoulli
polynomials. An individual multiple Gamma function also carries
polynomial coefficients in $x$. Theorem~\ref{bhr:thm:weighted-master}
supplies exactly those missing weights.

\subsection{Fixing the multiple Gamma normalization}

For an integer $r\ge1$, the unit-period Barnes zeta is initially
\begin{equation}\label{bhr:eq:barnes-zeta}
 \mathcal Z_r(u,x)=\sum_{n=0}^{\infty}
        \binom{n+r-1}{r-1}(n+x)^{-u},\qquad \Re u>r,\quad x>0.
\end{equation}
Define the explicit coefficient polynomials
\begin{equation}\label{bhr:eq:C-rk}
 C_{r,k}(x)=\frac1{(r-1)!}[X^k]
                      \prod_{h=1}^{r-1}(X+h-x),\qquad 0\le k<r.
\end{equation}
Expanding the binomial coefficient in powers of $n+x$ gives
\begin{equation}\label{bhr:eq:Barnes-Hurwitz}
 \mathcal Z_r(u,x)=\sum_{k=0}^{r-1}C_{r,k}(x)\zeta(u-k,x),\qquad
 L_r(x):=\partial_u\mathcal Z_r(0,x)
       =\sum_{k=0}^{r-1}C_{r,k}(x)\zeta'(-k,x).
\end{equation}
The first equality holds in the initial convergence domain and
then meromorphically. The value at zero is regular. This is the
classical Barnes/Hurwitz conversion, here written in a directly
computable polynomial form \cite{bhr:Adamchik,bhr:Vardi}.

To specify a normalization completely, put
\begin{equation}\label{bhr:eq:Q-r}
 Q_r(x)=\sum_{j=0}^{r-1}(-1)^{j+1}L_{r-j}(1)
                          \binom{x-1}{j},\qquad
 \log\Gamma_r(x)=L_r(x)+Q_r(x).
\end{equation}
We use the division convention for the multiple Gamma hierarchy:
$\Gamma_0(x)=x^{-1}$, $\Gamma_1=\Gamma$, and
$\Gamma_2=G^{-1}$. The zeta prescription
\eqref{bhr:eq:Barnes-Hurwitz}--\eqref{bhr:eq:Q-r}, not merely a functional
equation, defines the normalized functions in every higher rank.
Other standard sign conventions can be converted by inversion.

\begin{proposition}[Normalization and recurrence]\label{bhr:prop:normalization}
The functions in \eqref{bhr:eq:Q-r} satisfy
\begin{equation}\label{bhr:eq:Gamma-r-recurrence}
 \Gamma_r(1)=1,\qquad
 \Gamma_r(x+1)=\frac{\Gamma_r(x)}{\Gamma_{r-1}(x)}.
\end{equation}
Their normalization polynomials contain only rational numbers and
ordinary zeta derivatives at nonpositive integers.
\end{proposition}
\begin{proof}
Counting the multiplicities in \eqref{bhr:eq:barnes-zeta} gives
$\mathcal Z_r(u,x+1)-\mathcal Z_r(u,x)=-\mathcal Z_{r-1}(u,x)$, with
$\mathcal Z_0(u,x)=x^{-u}$. Hence
$L_r(x+1)-L_r(x)=-L_{r-1}(x)$, where $L_0=-\log x$.
The binomial difference formula gives
$Q_r(x+1)-Q_r(x)=-Q_{r-1}(x)$ for $r\ge2$.
For $r=1$, $Q_1=-\zeta'(0)=\tfrac12\log(2\pi)$ is
constant. Also $Q_r(1)=-L_r(1)$. These statements prove
\eqref{bhr:eq:Gamma-r-recurrence}, including the usual Gamma
recurrence at $r=1$. The last claim follows from
\eqref{bhr:eq:Barnes-Hurwitz} at $x=1$.
\end{proof}

The first three rows make every coefficient explicit. Here
$L=\log(2\pi)$:
\begin{align}
 L_1(x)&=\zeta'(0,x),& Q_1(x)&=L/2,\label{bhr:eq:rank1}\\
 L_2(x)&=\zeta'(-1,x)+(1-x)\zeta'(0,x),&
 Q_2(x)&=-\zeta'(-1)-\tfrac L2(x-1),\label{bhr:eq:rank2}\\
 L_3(x)&=\tfrac12\zeta'(-2,x)
       +\tfrac{3-2x}{2}\zeta'(-1,x)
       +\tfrac{(x-1)(x-2)}2\zeta'(0,x),\label{bhr:eq:rank3a}\\
 Q_3(x)&=-\tfrac12\bigl(\zeta'(-2)+\zeta'(-1)\bigr)
       +(x-1)\zeta'(-1)+\tfrac L2\binom{x-1}{2}.
       \label{bhr:eq:rank3b}
\end{align}
At rank two these equations are equivalent to the usual Vardi
identity
\begin{equation}\label{bhr:eq:Vardi}
 \log G(x)=\zeta'(-1)+(x-1)\log\Gamma(x)-\zeta'(-1,x).
\end{equation}
The familiar Barnes function has its standard asymptotic
normalization; $G(x+1)=\Gamma(x)G(x)$ and $G(1)=1$ alone
would allow multiplication by a nontrivial periodic factor.

\subsection{The all-rank resolvent theorem}

\begin{theorem}[Individual multiple Gamma transforms]\label{bhr:thm:all-Barnes}
For every integer $r\ge1$, polynomial $W$, and $z\notin\R$,
\begin{equation}\label{bhr:eq:all-Barnes}
 \boxed{\displaystyle
 \int_0^1 W(x)\log\Gamma_r(x)R_z(x)\,dx
  =-\sum_{k=0}^{r-1}
      \left.\partial_s I_{W C_{r,k}}(s,z)\right|_{s=k+1}
          +T_{WQ_r}(z).}
\end{equation}
Every term on the right is explicitly given by
\eqref{bhr:eq:weighted-master} and \eqref{bhr:eq:Tpoly}. It is a finite
combination of ordinary and depth-two polylogarithms, their first
order derivatives, Gamma factors, logarithms, and ordinary zeta
values and derivatives. All depth-two series needed at the stated
evaluation points are normally convergent.
\end{theorem}

\begin{proof}
Differentiate the ordinary integral in
Theorem~\ref{bhr:thm:weighted-master} at $s=k+1$:
\[
 \partial_s I_P(k+1,z)
       =-\int_0^1P(x)\zeta'(-k,x)R_z(x)\,dx.
\]
The shift formula at zero bounds the differentiated integrand by
a fixed polynomial in $|\log x|$ times $x$ when $k=0$, and
by an even more integrable power when $k>0$; its analytic tail
is bounded. Thus differentiation is justified locally uniformly.
Insert \eqref{bhr:eq:Barnes-Hurwitz} and \eqref{bhr:eq:Q-r}.
At every evaluation point $s=k+1\ge1$ and every polynomial
index $j\ge1$, the depth-two total order is $s+j\ge2$.
Lemma~\ref{bhr:lem:omega-normal} therefore applies with a margin,
including all the stated spectral derivatives. The sums over
$r,k,j$ are finite.
\end{proof}

This theorem answers the individual, unit-period multiple Gamma
part of source question R6. It does not assert a reduction for
arbitrary products of two multiple Gamma logarithms, nor for
arbitrary incommensurable Barnes periods.

Repeated poles are equally explicit:
\begin{equation}\label{bhr:eq:Barnes-repeated}
 \int_0^1\frac{W(x)\log\Gamma_r(x)}{(K(x)-z)^h}\,dx
   =\frac1{(h-1)!}\partial_z^{h-1}
      \int_0^1W(x)\log\Gamma_r(x)R_z(x)\,dx,
 \qquad h\ge1.
\end{equation}
All these integrals converge. Argument derivatives of
$\mathcal O_{\alpha,\beta}(q)$ can be taken directly in its
series, multiplying by powers of $n$. The formal identity
\[
 q\partial_q\mathcal O_{\alpha,\beta}
  =\mathcal O_{\alpha-1,\beta}-\mathcal O_{\alpha,\beta-1}
\]
must be kept as a paired difference when either term separately
leaves its convergence domain. The original differentiated series
still converges normally. This qualification prevents an artificial
divergence in implementations at total order two.

\subsection{A compact formula for the ordinary Barnes function}

Only one depth-two derivative is needed at rank two. Define
\begin{equation}\label{bhr:eq:Dq}
 \mathcal D(q)=
    \left.(\partial_\alpha+\partial_\beta)
                      \mathcal O_{\alpha,\beta}(q)\right|_{\alpha=\beta=1}
  =-\sum_{n,m\ge1}\frac{q^n\log(m(m+n))}{m(m+n)}.
\end{equation}
It is holomorphic for $|q|<1$ and vanishes at zero. In the next
formulas all polylogarithms have argument $q$. Set
\begin{align}
 \mathcal A(q)
  &=\frac{\mathcal D(q)-\Li_2^{[1]}(q)
       -(\gamma+L+1)\Li_2(q)-\Li_1(q)\Li_1^{[1]}(q)}{2\pi^2}
           -\frac14\Li_1(q),\label{bhr:eq:Aq}\\
 \mathcal S(q)
  &=\frac{\Li_1^{[1]}(q)-(\gamma+2L)\Li_1(q)
            -\tfrac12\bigl(\Li_2(q)+\Li_1(q)^2\bigr)}{2\pi},
        \label{bhr:eq:Sq}\\
 \mu_G&=2\zeta'(-1)-\frac1{12}-\frac L4.
       \label{bhr:eq:meanG}
\end{align}

\begin{theorem}[The individual $G$ resolvent]\label{bhr:thm:G-resolvent}
With the coordinates \eqref{bhr:eq:coordinates},
\begin{equation}\label{bhr:eq:G-resolvent}
 \boxed{\displaystyle
 \int_0^1\frac{\log G(x)}{\pi\cot(\pi x)-z}\,dx
    =-\frac{\mu_G}{a}
       +\frac c{2q}\bigl(\mathcal A(q)+\epsilon\ii\mathcal S(q)\bigr).}
\end{equation}
The quotient at $q=0$ is removable. The formula is holomorphic
on each open half-plane and uses the real logarithm of $G(x)$
on the integration interval.
\end{theorem}

\begin{proof}
We give a summation proof that also connects the formula to the
known Barnes Fourier expansion \cite[Theorems 2--3]{bhr:Mezo}.
Write
\[
 \log G(x)=\mu_G+\sum_{n\ge1}
                  \bigl(a_n\cos(2\pi nx)+b_n\sin(2\pi nx)\bigr).
\]
Its coefficients are
\begin{align}
 a_n&=\frac{\frac12\log n-(\gamma+L)-1}{2\pi^2n^2}
                -\frac1{4n}-\frac{T_n}{\pi^2},\label{bhr:eq:anG}\\
 b_n&=\frac{\frac1{2n}-\gamma-2L-\log n-H_n}{2\pi n},
          \label{bhr:eq:bnG}\\
 T_n&=\sum_{\substack{m\ge1\\m\ne n}}
                          \frac{\log m}{m^2-n^2}.
          \label{bhr:eq:Tn}
\end{align}
The $m$-sum in $T_n$ converges absolutely. These coefficients
can also be recovered from Theorem~\ref{bhr:thm:weighted-master}
by differentiating the $P=x-1$ formula at $s=1$, the $P=1$
formula at $s=2$, and using \eqref{bhr:eq:Vardi}. Thus their use
does not assume a new Barnes Fourier theorem.

To evaluate their generating functions, differentiate the
$j=1$ case of \eqref{bhr:eq:lattice-splitting} and the definition
of $\mathcal O_{1,s}$. Partial fractions give the exact identity
\begin{equation}\label{bhr:eq:T-generator}
 \sum_{n\ge1}q^nT_n
   =\frac12\Li_1(q)\Li_1^{[1]}(q)
       -\frac12\mathcal D(q)+\frac14\Li_2^{[1]}(q).
\end{equation}
For clarity, the scalar calculation behind it is
\[
 \frac1{m^2-n^2}=
     \frac1{2m}\left(\frac1{m-n}+\frac1{m+n}\right).
\]
The excluded term $m=n$ in the second sum is
$\log n/(2n^2)$ before multiplication by $1/2$; it supplies
the last term in \eqref{bhr:eq:T-generator}. This diagonal term
is necessary for the coefficient of $\Li_2^{[1]}$.

Next,
\begin{equation}\label{bhr:eq:harmonic-generator}
 \sum_{n\ge1}\frac{H_nq^n}{n}
      =\Li_2(q)+\frac12\Li_1(q)^2.
\end{equation}
It follows either by integrating
$\sum_{n\ge1}H_nq^n=\Li_1(q)/(1-q)$ or by the depth-two
stuffle identity. Equations \eqref{bhr:eq:anG}--\eqref{bhr:eq:harmonic-generator}
give $\sum a_nq^n=\mathcal A(q)$ and
$\sum b_nq^n=\mathcal S(q)$.

The function $\log G$ is integrable and square integrable:
near zero it equals $\log x+O(x)$, while it is analytic at one.
Pair its Fourier coefficients with the uniformly convergent
geometric series \eqref{bhr:eq:resolvent-fourier}. This yields
\eqref{bhr:eq:G-resolvent}. The constant term is
\eqref{bhr:eq:meanG}; for example integrate \eqref{bhr:eq:Vardi}, use
$\int_0^1\zeta'(-1,x)\,dx=0$ and
\[
 \int_0^1x\log\Gamma(x)\,dx
       =\frac L4+\zeta'(-1)-\frac1{12}.
\]
The latter follows by integrating the Hurwitz argument ladder.
All generating functions used above are $O(q)$, so the quotient
at zero is removable. Normal convergence proves holomorphy.
\end{proof}

\subsection{An ordinary-polylogarithm real component}

\begin{corollary}[A one-parameter exact real integral]\label{bhr:cor:G-real}
For $b>0$, let $q=(b-\pi)/(b+\pi)\in(-1,1)$. Then
\begin{equation}\label{bhr:eq:G-real}
 \boxed{\displaystyle
 \int_0^1\frac{K(x)\log G(x)}{K(x)^2+b^2}\,dx
 =\frac{\Li_1^{[1]}(q)-(\gamma+2L)\Li_1(q)
       -\frac12\bigl(\Li_2(q)+\Li_1(q)^2\bigr)}
               {2q(b+\pi)^2}.}
\end{equation}
At $b=\pi$ the right side means
$-(\gamma+2L+1/2)/(8\pi^2)$. In particular,
\begin{equation}\label{bhr:eq:G-threepi}
 \boxed{\displaystyle
 \int_0^1\frac{\pi\cot(\pi x)\log G(x)}
                      {\pi^2\cot^2(\pi x)+9\pi^2}\,dx
 =\frac{\Li_1^{[1]}(1/2)-(\gamma+2\log(2\pi))\log2
                    -\pi^2/24-(\log2)^2/4}{16\pi^2}.}
\end{equation}
\end{corollary}
\begin{proof}
At $z=\ii b$, $a=\ii(b+\pi)$ and
$c=-2\pi\ii/(b+\pi)^2$. The functions $\mathcal A,
\mathcal S$ are real at real $q$. Taking the real part of
\eqref{bhr:eq:G-resolvent} therefore keeps only
$\pi\mathcal S(q)/(q(b+\pi)^2)$.
Use $\Li_1(q)/q\to1$, $\Li_2(q)/q\to1$ and
$\Li_1^{[1]}(q)/q\to0$ at zero. For $b=3\pi$, use
$q=1/2$, $\Li_1(1/2)=\log2$ and
$\Li_2(1/2)=\pi^2/12-(\log2)^2/2$.
\end{proof}

The integral \eqref{bhr:eq:G-threepi} is approximately
$-0.02312768479852353748486907193849$.
The complete complex transform at $z=3\pi\ii$ is approximately
\[
 -0.02312768479852353748486907193849
   -0.05510836371472001971368271205301\,\ii.
\]
These decimals are diagnostics, not definitions or exact-value
certificates. The symmetry of $K$ under $x\mapsto1-x$ explains
why its real component isolates the simpler sine coefficients.

\subsection{Normalized antiderivatives remain explicit}

The standard balanced primitives are
\begin{equation}\label{bhr:eq:balanced-B}
 \mathscr B_k(x)=\frac{\zeta'(-k,x)+H_k\zeta(-k,x)}{k!},
 \qquad H_0=0.
\end{equation}
They satisfy $\mathscr B_0=\log\Gamma-L/2$ and
$\mathscr B_k'=\mathscr B_{k-1}$ for $k\ge1$
\cite{bhr:EMIntegrals,bhr:EMGeneralized}. For a polynomial $P$, set
\begin{equation}\label{bhr:eq:weighted-antiderivative}
 \mathscr A_{P,k}(x)=
    \sum_{j=0}^{\deg P}(-1)^jP^{(j)}(x)\mathscr B_{k+j+1}(x).
\end{equation}
Termwise differentiation telescopes, so
$\mathscr A_{P,k}'=P\mathscr B_k$.
Consequently a primitive of $P(x)\zeta'(-k,x)$, normalized
to vanish at $x=1$, is
\begin{align}\label{bhr:eq:normalized-zeta-primitive}
 k!\bigl(\mathscr A_{P,k}(x)-\mathscr A_{P,k}(1)\bigr)
    -H_k\int_1^xP(t)\zeta(-k,t)\,dt.
\end{align}
The remaining integral is polynomial because
$\zeta(-k,t)=-B_{k+1}(t)/(k+1)$.
Apply \eqref{bhr:eq:normalized-zeta-primitive} to the finitely many
terms in \eqref{bhr:eq:Barnes-Hurwitz}, and integrate $Q_r$ as a
polynomial. This gives a fully normalized primitive of every
$\log\Gamma_r$. Repetition gives every iterated primitive,
without unspecified additive constants. Their resolvent
transforms again follow from Theorem~\ref{bhr:thm:all-Barnes} and
Theorem~\ref{bhr:thm:weighted-master}. The primitive operation is
classical integration by parts; the point here is its compatibility
with the explicit individual-function normalization.

% END sections/03_barnes.tex


% BEGIN sections/04_complex_powers.tex
\section{Complex cotangent powers and their resonant amplitudes}
\label{bhr:cxp:sec:complex}

Question R4 of the incoming \emph{Harmonic Parity and Resolvent
Identities} report \cite{bhr:Parity} asks for a joint generator for nonintegral powers of
the cotangent resolvent, with the moving endpoint amplitudes retained
before expansion.  We prove such a generator and its finite-part
conversion.  The primary result below covers every complex exponent
with positive real part and every positive integral crossing.  A new
Dirichlet-series coordinate is used explicitly; no finite reduction to
ordinary polylogarithms is asserted for an arbitrary complex exponent.
Nevertheless, its exponent jets have a finite multiple-polylogarithmic
description, and one special resolvent point has a Gamma evaluation at
every positive integral spectral order.

\subsection{Branches and a binomial Dirichlet transform}

Put
\[
 K(x)=\pi\cot(\pi x),\qquad R_z(x)=(K(x)-z)^{-1},
 \qquad z\in\mathbb C\setminus\mathbb R.
\]
For the chosen half-plane define
\[
 \epsilon=\operatorname{sgn}\Im z,\quad
 a=z+\epsilon\pi i,\quad A=-a^{-1},\quad
 q=\frac{z-\epsilon\pi i}{z+\epsilon\pi i},\quad
 \sigma=-2\epsilon\pi i.
\]
Then \(|q|<1\).  We take the continuous logarithm of \(K(x)-z\)
whose imaginary part lies in \((-\pi,0)\) if \(\epsilon=1\) and
in \((0,\pi)\) if \(\epsilon=-1\).  Consequently
\[
 R_z(x)^\lambda=\exp\{-\lambda\Log(K(x)-z)\}.
\]
Also \(\Log\sigma=\log(2\pi)-\epsilon\pi i/2\), and
\(\Log A\) is the principal logarithm.

With \(w=e^{2\epsilon\pi i x}\), elementary algebra gives
\begin{equation}\label{bhr:cxp:eq:disk}
 R_z(x)^\lambda
 =A^\lambda(1-w)^\lambda(1-qw)^{-\lambda}.
\end{equation}
Here the two logarithms on the right are their holomorphic
determinations in \(|w|<1\), continued to \(w\ne1\) on its boundary.
To check the constant in this identity, both sides divided by
\(x^\lambda\) tend to one as \(x\downarrow0\).  This fixes the phase
at the other endpoint as well.

Define
\begin{align}
 h_n(\lambda,q)
 &=[w^n](1-w)^\lambda(1-qw)^{-\lambda}\notag\\
 &=\sum_{k=0}^n
   \frac{(-\lambda)_k(\lambda)_{n-k}}{k!(n-k)!}\,q^{n-k},
 \label{bhr:cxp:eq:hn}\\
 \mathcal L_\lambda(u,q)
 &=\sum_{n=1}^{\infty}\frac{h_n(\lambda,q)}{n^u}.
 \label{bhr:cxp:eq:L}
\end{align}
The Pochhammer symbols have their polynomial meaning, including at
integral parameters.  The initial Dirichlet series converges normally
where \(\Re(u+\lambda)>0\), locally in \(|q|<1\).  One way to see
the needed coefficient estimate is to convolve the binomial
coefficients of \((1-w)^\lambda\) with the exponentially decaying
coefficients of \((1-qw)^{-\lambda}\); on compact parameter sets this
gives \(h_n=O(n^{-\Re\lambda-1})\), with a harmless uniform weakening
of the exponent at the boundary of a compact set.

\begin{theorem}[Joint continuation of the binomial transform]
\label{bhr:cxp:thm:L}
The function \(\mathcal L_\lambda(u,q)\) continues meromorphically
in \((u,\lambda)\in\mathbb C^2\), holomorphically in \(|q|<1\).
Its possible polar hypersurfaces are
\[
 u+\lambda=-j,\qquad j=0,1,2,\ldots.
\]
Away from them,
\[
 \mathcal L_\lambda(0,q)=-1,\qquad
 \mathcal L_\lambda(-N,q)=0\quad(N\ge1).
\]
These equalities concern the meromorphic restriction at the indicated
spectral order.  At an intersection with a polar hypersurface,
substitution and continuation need not commute.
\end{theorem}

\begin{proof}
For \(\Re u>0\) and \(\Re(u+\lambda)>0\), termwise integration
gives
\begin{equation}\label{bhr:cxp:eq:mellin}
 \Gamma(u)\mathcal L_\lambda(u,q)
 =\int_0^\infty t^{u-1}
 \left[
 \left(\frac{1-e^{-t}}{1-qe^{-t}}\right)^\lambda-1
 \right]\,dt.
\end{equation}
The positive real logarithm of \(1-e^{-t}\), and the holomorphic
logarithm of \(1-qe^{-t}\), specify its powers.  Set
\[
 \left(\frac{1-e^{-t}}{1-qe^{-t}}\right)^\lambda
   =t^\lambda\sum_{j\ge0}b_j(\lambda,q)t^j.
\]
The coefficients are entire in \(\lambda\), holomorphic in \(q\).
For any fixed \(N\), split the integral at one and subtract the
first \(N\) terms of this local expansion.  The result is
\begin{align}
 \Gamma(u)\mathcal L_\lambda(u,q)
 ={}&-\frac1u+
   \sum_{j=0}^{N-1}\frac{b_j(\lambda,q)}{u+\lambda+j}
   +E_N(u,\lambda,q),\label{bhr:cxp:eq:subtraction}
\end{align}
where \(E_N\) is holomorphic when \(\Re(u+\lambda)>-N\).
The tail at infinity is exponentially convergent on every compact
parameter set.  Increasing \(N\) proves the continuation.  Division
by \(\Gamma(u)\) removes the isolated pole \(-1/u\), with value
\(-1\) at \(u=0\), and gives zero at the other nonpositive
integers whenever no moving denominator vanishes.
\end{proof}

\subsection{The completed spectral generator}

Use the conventional normalization
\[
 \zeta(1-s,x)=-\frac1s+
    \sum_{m\ge0}\gamma_m(x)\frac{s^m}{m!},
 \qquad \gamma_0(x)=-\psi(x).
\]
For an integer \(p\ge0\), let
\[
 A_p(s)=\prod_{k=1}^p(s-k),\qquad A_0(s)=1.
\]
Argument differentiation satisfies
\[
 \partial_x^p\zeta(1-s,x)
   =A_p(s)\zeta(1+p-s,x).
\]
The notation \(\gamma_m^{(p)}\) below always means argument
differentiation.

\begin{theorem}[Complex-power cotangent generator]
\label{bhr:cxp:thm:generator}
In the ordinary-integral region
\(\Re\lambda>-1,\ \Re(s+\lambda)>p\), away from the Hurwitz
pole when \(p=0\),
\begin{equation}\label{bhr:cxp:eq:J}
 \mathcal J_p(s,\lambda,z):=
 \int_0^1\partial_x^p\zeta(1-s,x)R_z(x)^\lambda\,dx
 =A^\lambda\Gamma(s)\sigma^{p-s}
       \mathcal L_\lambda(s-p,q).
\end{equation}
This identity defines a joint meromorphic continuation to all
\((s,\lambda)\).  Besides the pole
\(-A^\lambda/s\) when \(p=0\), its only possible polar
hypersurfaces are
\[
 s+\lambda=p-j,\qquad j=0,1,\ldots.
\]
Define the completed family
\begin{equation}\label{bhr:cxp:eq:F}
 \mathcal F_p(s,\lambda,z)=
 \mathcal J_p(s,\lambda,z)
       +\mathbf1_{p=0}\frac{A^\lambda}{s}.
\end{equation}
For \(\Re\lambda>p\) the ordinary absolutely convergent moments
\begin{equation}\label{bhr:cxp:eq:M}
 \mathcal M_{m,p}(\lambda,z)=
 \int_0^1\gamma_m^{(p)}(x)R_z(x)^\lambda\,dx
 =m![s^m]\mathcal F_p(s,\lambda,z)
\end{equation}
therefore have meromorphic continuations in \(\lambda\), with
possible poles only at the integers \(\lambda\le p\), of order
at most \(m+1\).
\end{theorem}

\begin{proof}
First suppose \(\Re\lambda>0\) and \(\Re s>p+1\).
The Fourier series of the kernel in \eqref{bhr:cxp:eq:disk}
is absolutely convergent.  The classical Hurwitz Fourier identity
\cite[25.11.9]{bhr:DLMFZ} gives
\[
 \int_0^1\zeta(1-s,x)\,dx=0,\qquad
 \int_0^1\partial_x^p\zeta(1-s,x)
           e^{2\epsilon\pi i n x}\,dx
 =\Gamma(s)\sigma^{p-s}n^{p-s}.
\]
The constant kernel mode vanishes, and summing the other modes
proves \eqref{bhr:cxp:eq:J}.  The endpoint decomposition
\[
 \partial_x^p\zeta(1-s,x)
 =A_p(s)x^{s-p-1}
       +\partial_x^p\zeta(1-s,1+x)
\]
shows normal integrability on the larger region stated in the
theorem, so the identity extends there.

In \eqref{bhr:cxp:eq:subtraction}, put \(u=s-p\) and use
\(\Gamma(s)/\Gamma(s-p)=A_p(s)\).  The moving denominators
give exactly the indicated hypersurfaces.  The isolated term
\(-1/(s-p)\) is cancelled by \(A_p(s)\) for \(p\ge1\); for
\(p=0\) it gives residue \(-A^\lambda\).  This proves the
claimed completion.  Finally, the mean of the kernel is \(A^\lambda\)
for \(\Re\lambda>0\), by its constant Fourier coefficient.
Termwise Stieltjes expansion is valid when \(\Re\lambda>p\),
including all powers of the endpoint logarithm.  This proves
\eqref{bhr:cxp:eq:M} and its continuation.
\end{proof}

\begin{corollary}[Lowest Stieltjes index and shift identities]
\label{bhr:cxp:cor:shifts}
Away from the moving poles,
\begin{align*}
 \mathcal M_{0,0}(\lambda,z)
 &=A^\lambda\left(
   \gamma+\Log\sigma+
         \left.\partial_u\mathcal L_\lambda(u,q)\right|_{u=0}\right),\\
 \mathcal M_{0,p}(\lambda,z)
 &=A^\lambda\sigma^p
         \left.\partial_u\mathcal L_\lambda(u,q)\right|_{u=-p}
                                      \qquad(p\ge1).
\end{align*}
The continued moments satisfy the meromorphic identities
\begin{align}
 \partial_z\mathcal M_{m,p}(\lambda,z)
   &=\lambda\mathcal M_{m,p}(\lambda+1,z),\label{bhr:cxp:eq:zshift}\\
 \mathcal M_{m,p+1}(\lambda,z)
   &=-\lambda\bigl(
       \mathcal M_{m,p}(\lambda-1,z)
       +2z\mathcal M_{m,p}(\lambda,z)\notag\\
   &\hspace{39mm}
       +(z^2+\pi^2)\mathcal M_{m,p}(\lambda+1,z)\bigr).
       \label{bhr:cxp:eq:pshift}
\end{align}
\end{corollary}
\begin{proof}
The first two formulas follow by expanding \(\Gamma(s)\sigma^{-s}\)
and using Theorem~\ref{bhr:cxp:thm:L} at \(0,-p\).
The first shift follows from
\(\partial_zR_z^\lambda=\lambda R_z^{\lambda+1}\).
For the second, integrate by parts where \(\Re\lambda>p+1\),
so the boundary terms vanish, and use
\[
 R_z'=1+2zR_z+(z^2+\pi^2)R_z^2.
\]
Continuation proves the meromorphic statements.  At a resonant
integer, a pointwise finite-part version must use the conversion
in the next theorem before substituting.
\end{proof}

\subsection{All positive resonances and the two finite parts}

Finite parts in the next theorem use the original \(x\)-coordinate:
they retain the constant coefficient in the cutoff integral
\(\int_\eta^{1-\eta}\), after removing every nonintegrable power
and logarithm.  For a complex exponent of positive real part this
definition is unambiguous from the local power-log expansion.
It agrees with the ordinary integral whenever that integral converges.

The complete local amplitude is elementary.  Write
\begin{equation}\label{bhr:cxp:eq:d}
 D_j(\lambda,z)
 =[x^j]\bigl(\pi x\cot(\pi x)-zx\bigr)^{-\lambda}.
\end{equation}
Thus \(R_z(x)^\lambda
=x^\lambda\sum_{j\ge0}D_j(\lambda,z)x^j\) near zero, with
\[
 D_0=1,\qquad D_1=\lambda z,\qquad
 D_2=\frac{\lambda\pi^2}{3}
       +\frac{\lambda(\lambda+1)z^2}{2}.
\]
Every \(D_j\) is a polynomial in \(\lambda,z,\pi^2\) with rational
coefficients, and has degree at most \(j\) in \(\lambda\).

\begin{theorem}[Resonant amplitudes and finite-part conversion]
\label{bhr:cxp:thm:contacts}
For a positive integer \(r\le p\), put \(j=p-r\).
In a neighborhood of \((s,\lambda)=(0,r)\),
\begin{equation}\label{bhr:cxp:eq:amplitude}
 \mathcal F_p(s,\lambda,z)
  =\frac{A_p(s)D_j(\lambda,z)}{s+\lambda-r}
      +H_{p,r}(s,\lambda,z),
\end{equation}
where \(H_{p,r}\) is jointly holomorphic.  Hence, for every
positive integer \(r\),
\begin{align}
 &\operatorname{FP}_x\int_0^1
       \gamma_m^{(p)}(x)R_z(x)^r\,dx\notag\\
 &\quad =
 m![s^m]\left\{
   \mathcal F_p(s,r,z)
   -\mathbf1_{r\le p}\frac{A_p(s)D_{p-r}(r,z)}s
           \right\}. \label{bhr:cxp:eq:fixed-r}
\end{align}

There is also an explicit comparison with continuation in the
kernel exponent.  Write \(A_p(s)=\sum_{h=0}^p a_{p,h}s^h\).
For \(1\le r\le p\), define
\begin{equation}\label{bhr:cxp:eq:Delta}
 \Delta_{m,p,r}(z)=m!
 \sum_{h=0}^{\min(p,m)}
 \frac{(-1)^{m-h}a_{p,h}}{(m-h+1)!}\,
 \left.\partial_\lambda^{\,m-h+1}
                 D_{p-r}(\lambda,z)\right|_{\lambda=r}.
\end{equation}
Then
\begin{equation}\label{bhr:cxp:eq:lambdaCT}
 \operatorname{FP}_x\int_0^1
   \gamma_m^{(p)}(x)R_z(x)^r\,dx
 =\operatorname{CT}_{\lambda=r}
      \mathcal M_{m,p}(\lambda,z)-\Delta_{m,p,r}(z).
\end{equation}
For \(\Re\lambda>0\) away from the positive integers
\(1,\ldots,p\), the meromorphic value of
\(\mathcal M_{m,p}\) equals the pointwise \(x\)-finite part.
\end{theorem}

\begin{proof}
Subtract a sufficiently long full local Taylor jet, so that the
remaining endpoint integral converges normally near \((0,r)\).
Among the explicitly integrated Taylor terms, the only one whose
exponent crosses \(-1\) there is
\[
 A_p(s)D_{p-r}(\lambda,z)
                 x^{s+\lambda-r-1}.
\]
Use a smooth cutoff on a small interval next to zero.
The continued integrals of all other Taylor terms, and the normally
convergent remainder, are holomorphic near \((0,r)\).
The distinguished term has meromorphic integral with polar part
\(A_p(s)D_{p-r}(\lambda,z)/(s+\lambda-r)\).
Replacing the cutoff by one only changes a holomorphic function;
this proves \eqref{bhr:cxp:eq:amplitude} with the indicated full
amplitude, not just its value on the polar hyperplane.

At fixed \(\lambda=r\), the analytically continued integral of
the resonant local term is \(A_p(s)D_{p-r}(r,z)/s\).
Every coefficient of \(x^{s-1}\) is a power of \(\log x\)
divided by \(x\), whose finite part on \((0,1)\) is zero.
Deleting the entire amplitude before expansion proves
\eqref{bhr:cxp:eq:fixed-r}.  Terms on the cutoff complement
cancel the artificial splitting point.

For the second formula, put \(\delta=\lambda-r\) and take
the \(s^m\) coefficient of the polar term in
\eqref{bhr:cxp:eq:amplitude}.  It is
\[
 m!\sum_{h=0}^{\min(p,m)}
  (-1)^{m-h}a_{p,h}
       \frac{D_{p-r}(r+\delta,z)}{\delta^{m-h+1}}.
\]
Its constant Laurent coefficient is exactly
\eqref{bhr:cxp:eq:Delta}.  The holomorphic remainder contributes
the same value along either order of expansion.  This proves
\eqref{bhr:cxp:eq:lambdaCT}.  Away from a crossing, endpoint
subtractions have nonzero exponents, and coefficient extraction
commutes with finite-part integration.
\end{proof}

The theorem supplies the full principal part as well: expand the
polynomial \(D_{p-r}(r+\delta,z)\) in the last display and retain
negative powers of \(\delta\).  Some potential poles cancel.
In particular, \(\Delta_{m,p,p}=0\), since \(D_0=1\);
a collision does not automatically imply a nonzero conversion term.

\paragraph{A nonzero conversion for the simple resolvent.}
The simplest example is \(m=0,p=2,r=1\).  Since
\(A_2(s)=s^2-3s+2\) and \(D_1(\lambda,z)=\lambda z\),
\[
 \Delta_{0,2,1}(z)=2z.
\]
At \(z=i\pi\), set \(\ell=\log(2\pi)-i\pi/2\).
The finite Fourier series
\(R_{i\pi}(x)=(e^{2\pi ix}-1)/(2\pi i)\) gives
\[
 \operatorname{FP}_x\int_0^1\gamma_0''(x)R_{i\pi}(x)\,dx
       =2\pi i\left(\frac32-\gamma-\ell\right).
\]
It follows that
\begin{equation}\label{bhr:cxp:eq:noncommute}
 \operatorname{CT}_{\lambda=1}
       \mathcal M_{0,2}(\lambda,i\pi)
       =2\pi i\left(\frac52-\gamma-\ell\right).
\end{equation}
The two constants differ by \(2\pi i\).  This is a concrete
obstruction to identifying exponent continuation with the
specified \(x\)-finite part.  In contrast the case
\(m=0,p=r=1\) has no exponent-axis correction.

\paragraph{Exact recovery of integral powers.}
For a positive integer \(r\) and \(q\ne0\), partial fractions give
\[
 h_n(r,q)=q^{n-r}
   \sum_{j=1}^r\binom rj(q-1)^j
                    \binom{n+j-1}{j-1}\qquad(n\ge1).
\]
The binomial on the right is a polynomial of degree \(j-1\) in
\(n\).  Consequently \(\mathcal L_r(u,q)\) is a finite linear
combination of the ordinary \(\operatorname{Li}_{u-k}(q)\),
\(0\le k<r\), with rational coefficients in \(q\).
At \(q=0\) the removable limiting expression is
\[
 \mathcal L_r(u,0)
       =\sum_{n=1}^r\frac{(-1)^n\binom rn}{n^u}.
\]
Thus the completed generator recovers precisely the ordinary
polylogarithm class of the previous repeated-resolvent theorem.

\subsection{Gamma closure at a special point}

For \(z=\epsilon i\pi\) one has \(q=0\), and
\[
 \mathcal L_\lambda(u,0)
       =\sum_{n\ge1}\frac{(-\lambda)_n}{n!\,n^u}.
\]
At every positive integer \(u\), this new coordinate reduces to
generalized harmonic numbers at a complex argument.  Define
\[
 H_\lambda^{(1)}=\psi(1+\lambda)+\gamma,\qquad
 H_\lambda^{(r)}=\zeta(r)-\zeta(r,1+\lambda)\quad(r\ge2),
\]
initially for \(\Re\lambda>-1\) and then meromorphically.

\begin{theorem}[Every integral spectral order at \(q=0\)]
\label{bhr:cxp:thm:gamma}
For \(N\ge1\),
\begin{equation}\label{bhr:cxp:eq:Bell}
 \mathcal L_\lambda(N,0)
 =-\frac1{N!}Y_N\bigl(
     H_\lambda^{(1)},1!H_\lambda^{(2)},\ldots,
                          (N-1)!H_\lambda^{(N)}\bigr),
\end{equation}
where \(Y_N\) is the complete exponential Bell polynomial.
Equivalently, near \(v=0\),
\begin{equation}\label{bhr:cxp:eq:beta-generator}
 \sum_{N\ge1}\mathcal L_\lambda(N,0)v^{N-1}
 =\frac1v\left\{
 1-\frac{\Gamma(1-v)\Gamma(1+\lambda)}
                 {\Gamma(1+\lambda-v)}\right\}.
\end{equation}
In particular,
\[
 \mathcal L_\lambda(1,0)=-H_\lambda^{(1)},\qquad
 \mathcal L_\lambda(2,0)
 =-\frac{(H_\lambda^{(1)})^2+H_\lambda^{(2)}}2.
\]
\end{theorem}

\begin{proof}
For \(\Re\lambda>-1\) and \(\Re t>0\), the beta integral gives
\[
 \sum_{n\ge1}\frac{(-\lambda)_n}{n!(n+t)}
  =\int_0^1 y^{t-1}\bigl((1-y)^\lambda-1\bigr)\,dy
  =\frac{\Gamma(t)\Gamma(1+\lambda)}
           {\Gamma(1+\lambda+t)}-\frac1t.
\]
Expand at \(t=-v\).  The left side generates the positive
integral spectral orders.  On the right,
\[
 \log\frac{\Gamma(1-v)\Gamma(1+\lambda)}
              {\Gamma(1+\lambda-v)}
   =\sum_{r\ge1}\frac{H_\lambda^{(r)}}r\,v^r.
\]
The defining Bell-polynomial generating function proves both
identities.  Meromorphic continuation removes the initial
restrictions.  This is a beta-integral specialization, not a claim
of priority for the classical binomial-harmonic identity.
\end{proof}

\subsection{Exponent jets and colored multiple polylogarithms}

There is also an exact finite description of every exponent jet at
\(\lambda=0\).  For \(q\ne0,\ |q|<1\), let
\(\mathcal A=\{1,q\}\), with signs \(\nu(1)=-1,\nu(q)=1\).
The multiple polylogarithm uses decreasing indices:
\[
 \operatorname{Li}_{s_1,\ldots,s_d}(z_1,\ldots,z_d)
 =\sum_{n_1>\cdots>n_d\ge1}
       \frac{z_1^{n_1}\cdots z_d^{n_d}}
            {n_1^{s_1}\cdots n_d^{s_d}}.
\]

\begin{proposition}[All exponent jets at zero]
\label{bhr:cxp:prop:jets}
For \(d\ge1\) and \(\Re u>0\),
\begin{equation}\label{bhr:cxp:eq:MPL}
 \left.\partial_\lambda^d\mathcal L_\lambda(u,q)
                                      \right|_{\lambda=0}
 =d!\sum_{a_1,\ldots,a_d\in\mathcal A}
 \left(\prod_{j=1}^d\nu(a_j)\right)
 \operatorname{Li}_{u+1,\{1\}^{d-1}}
       \left(a_1,\frac{a_2}{a_1},\ldots,
                                \frac{a_d}{a_{d-1}}\right).
\end{equation}
All displayed nested sums converge absolutely.  Their cumulative
color products are \(a_j\in\{1,q\}\), so the appearance of
\(q^{-1}\) as an individual color causes no convergence problem.
For nonintegral \(u\), these are nested Dirichlet sums with a complex first index; at positive integral \(u\), they are classical colored multiple polylogarithms.
At \(q=0\), the continuous limiting identity is
\[
 \left.\partial_\lambda^d\mathcal L_\lambda(u,0)
                                      \right|_{\lambda=0}
       =(-1)^d d!\,\zeta(u+1,\{1\}^{d-1}).
\]
\end{proposition}

\begin{proof}
Put \(F(w)=\operatorname{Li}_1(qw)-\operatorname{Li}_1(w)\).
The \(d\)-th exponent derivative of the disk kernel at zero is
\(F(w)^d\).  Repeated integration of
\[
 F'(w)=\frac{q}{1-qw}-\frac1{1-w}
\]
proves the shuffle identity
\[
 \frac{F(w)^d}{d!}
 =\sum_{a_1,\ldots,a_d\in\mathcal A}
   \left(\prod_j\nu(a_j)\right)
 \operatorname{Li}_{\{1\}^d}
     \left(a_1w,\frac{a_2}{a_1},\ldots,
                                   \frac{a_d}{a_{d-1}}\right).
\]
Multiplying the coefficient of \(w^n\) by \(n^{-u}\) raises
the first index from one to \(u+1\).  Absolute convergence for
\(\Re u>0\) justifies evaluation at \(w=1\) and all coefficient
interchanges.  If \(q=0\), only the letter one survives.
\end{proof}

The first jet is the simple identity
\[
 \left.\partial_\lambda\mathcal L_\lambda(u,q)\right|_0
       =\operatorname{Li}_{u+1}(q)-\zeta(u+1).
\]
Higher jets explain how generalized harmonic sums and increasing
multiple-polylogarithm depth enter complex powers.  The finite
description is proved for each fixed jet; it does not imply finite
ordinary-polylogarithm closure after summing all jets.

\paragraph{Verification.}
The accompanying program independently compares the spectral
generator with ordinary \(x\)-quadrature at nonintegral complex
parameters, checks the Gamma specialization against Mellin
quadrature, and compares a nearby-exponent Laurent extrapolation
with \eqref{bhr:cxp:eq:noncommute}.  Its recorded numerical residuals
are diagnostics.  The continuation, pole classification, and
all-order conversion rest on the local subtractions above.

% END sections/04_complex_powers.tex


% BEGIN sections/05_quartic_harmonic.tex
\section{The fourth digamma moment and harmonic Laurent data}
\label{bhr:quartic-harmonic:sec}

The incoming \emph{Cubic Tornheim Derivatives and Harmonic Stieltjes Constants}
report, Section~7,
Question~6 \cite{bhr:Cubic}, explicitly asks for a power-to-harmonic formula at the fourth
power. The following theorem answers that question. It is separate from
the quartic \emph{directional} Tornheim layer, which has already been
worked out in later incoming reports. No claim of worldwide novelty or
finite reduction to ordinary zeta constants is intended.

Write $H_n^{(r)}=\sum_{j=1}^n j^{-r}$, $H_n=H_n^{(1)}$, and
\[
 H_n(1,1)=\sum_{n\ge j>k\ge1}\frac1{jk}
         =\frac{H_n^2-H_n^{(2)}}2.
\]
The strict harmonic Dirichlet series used here are
\[
 Z_2(s,1)=\sum_{n=1}^{\infty}\frac{H_{n-1}}{n^s},\qquad
 Z_3(s,1,1)=\sum_{n=1}^{\infty}\frac{H_{n-1}(1,1)}{n^s},
 \qquad \Re s>1.
\]
They have the local expansions
\begin{align}
 Z_2(1+u,1)&=u^{-2}+\gamma u^{-1}
       +\frac{\gamma^2-\zeta(2)}2+\eta u+O(u^2),
       \label{bhr:quartic-harmonic:eta}\\
 Z_3(1+u,1,1)&=u^{-3}+\gamma u^{-2}
       +\frac{\gamma^2-\zeta(2)}{2u}
       +\frac{\gamma^3-3\gamma\zeta(2)+2\zeta(3)}6
       +\beta_1u+O(u^2).
       \label{bhr:quartic-harmonic:beta}
\end{align}
In particular $\beta_1$ is a completely specified strict depth-three
harmonic Laurent coefficient; it is not a coefficient of the inclusive
harmonic series. The constant term is also fixed by the Mellin lemma below.
The pole terms in the second expansion follow directly
from $H_{n-1}=\log n+\gamma+O(n^{-1})$ and
$H_{n-1}^{(2)}=\zeta(2)+O(n^{-1})$, by subtraction of zeta derivatives.
The convention for the ordinary Stieltjes constants is
\[
 \zeta(1+u)=u^{-1}+\sum_{j\ge0}\frac{(-1)^j\gamma_j}{j!}u^j.
\]

\begin{theorem}[A quartic power-to-harmonic identity]
\label{bhr:quartic-harmonic:main}
Use the unit-coordinate Hadamard finite part at the origin and put
$Q_4=\operatorname{FP}\int_0^1\psi(x)^4\,dx$. Then
\begin{align}
 Q_4={}&18\zeta(4)+12\zeta(3)-12\gamma\zeta(2)+4\zeta'(3)
       +12(\gamma^2-\zeta(2))\gamma_1
       +24\gamma\eta-24\beta_1.
       \label{bhr:quartic-harmonic:main-eq}
\end{align}
An equivalent absolutely convergent arithmetic formula is
\begin{align}
 Q_4={}&18\zeta(4)+12\zeta(3)-12\gamma\zeta(2)+4\zeta'(3)
          +12\gamma_3-24\zeta(2)\gamma_1\notag\\
 &+12\sum_{n=1}^{\infty}
       \frac{\bigl(\psi(n)^2+\psi'(n)-\log^2 n\bigr)\log n}{n}.
       \label{bhr:quartic-harmonic:arithmetic}
\end{align}
If $Q_3=\operatorname{FP}\int_0^1\psi(x)^3\,dx$, the already proved
cubic bridge \cite{bhr:Cubic} $Q_3=3\zeta(2)-6\eta-6\gamma\gamma_1$ eliminates $\eta$:
\begin{equation}
 Q_4+4\gamma Q_3
   =18\zeta(4)+12\zeta(3)+4\zeta'(3)
       -12(\gamma^2+\zeta(2))\gamma_1-24\beta_1.
       \label{bhr:quartic-harmonic:combined}
\end{equation}
\end{theorem}

\subsection{A normalized Mittag--Leffler expansion}

For $n\ge0$ set
\begin{align*}
 b_n&=H_n-\gamma,& c_n&=\zeta(2)+H_n^{(2)},&
 e_n&=H_n^{(3)}-\zeta(3),\\
 a_n&=c_n-b_n^2,& d_n&=6b_n^2-4c_n,&
 r_n&=-4b_n^3+12b_nc_n-4e_n.
\end{align*}
Let
\[
 C_* =\gamma^4-12\gamma^2\zeta(2)+12\gamma\zeta(3)+11\zeta(4),
 \qquad
 P_n(z)=\frac1{(n+z)^4}-\frac{4b_n}{(n+z)^3}
          +\frac{d_n}{(n+z)^2}+\frac{r_n}{n+z}.
\]

\begin{lemma}[Quartic principal parts with no unknown entire term]
\label{bhr:quartic-harmonic:ML}
For $z\notin\{0,-1,-2,\ldots\}$,
\begin{equation}
 \psi(z)^4=P_0(z)+C_*+
     \sum_{n=1}^{\infty}\bigl(P_n(z)-P_n(0)\bigr).
     \label{bhr:quartic-harmonic:ML-eq}
\end{equation}
The series converges normally on every compact set avoiding the poles.
\end{lemma}

\begin{proof}
The digamma recurrence and its convergent expansion at zero \cite{bhr:DLMFG} give
\[
 \psi(-n+\varepsilon)
   =-\varepsilon^{-1}+b_n+c_n\varepsilon+e_n\varepsilon^2
       +(\zeta(4)+H_n^{(4)})\varepsilon^3+O(\varepsilon^4).
\]
Raising this expression to the fourth power gives exactly $P_n$ as
its principal part. The constant coefficient at $n=0$ is
\[
 \gamma^4-12\gamma^2\zeta(2)+6\zeta(2)^2
                +12\gamma\zeta(3)-4\zeta(4)=C_*.
\]
This establishes the poles but does not on its own establish the entire
remainder. To determine that remainder, let $f(z)=\psi(z)^4-P_0(z)$.
For fixed $z$, apply the residue theorem to
$z f(w)/(w(w-z))$ on $|w|=N+\tfrac12$.
The standard digamma asymptotic in the right half-plane, and
$\psi(w)=\psi(1-w)-\pi\cot(\pi w)$ in the left half-plane, give
$f(w)=O(\log^4 N)$ uniformly on these circles. The cotangent is
uniformly bounded there: away from the real axis this follows from its
exponential formula, while near the real axis these half-integer circles
remain a fixed positive distance from the integer poles. Thus the
contour integral is $O(\log^4N/N)$ and tends to zero. The residues are
$f(z)$, $-f(0)$, and $-(P_n(z)-P_n(0))$, proving the formula.
Finally $b_n=O(\log n)$, $d_n=O(\log^2n)$ and
$r_n=O(\log^3n)$, and the simple-pole difference is $O(n^{-2})$
on each fixed compact set. These estimates prove normal convergence.
\end{proof}

\subsection{The exact summation-by-parts cancellation}

For $m\ge2$ and $n\ge1$,
\[
 \int_0^1\bigl((n+x)^{-m}-n^{-m}\bigr)\,dx
     =\frac{n^{1-m}-(n+1)^{1-m}}{m-1}-n^{-m}.
\]
Integrating \eqref{bhr:quartic-harmonic:ML-eq}, and retaining the finite
parts of its three higher-order poles at zero, yields
\begin{equation}
 Q_4=K+\sum_{n=1}^{\infty}r_n
      \left(\log\left(1+\frac1n\right)-\frac1n\right),
      \label{bhr:quartic-harmonic:K-eq}
\end{equation}
where
\begin{equation}
 K=\gamma^4-18\gamma^2\zeta(2)+32\gamma\zeta(3)
       +\frac{13}{2}\zeta(4)+12\zeta(3)-12\gamma\zeta(2).
       \label{bhr:quartic-harmonic:K}
\end{equation}
For clarity, the ordinary Euler sums needed in this simplification are
\begin{equation}
 \sum_{n\ge1}\frac{H_n}{n^2}=2\zeta(3),\qquad
 \sum_{n\ge1}\frac{H_n}{n^3}=\frac54\zeta(4),\qquad
 \sum_{n\ge1}\frac{H_n^{(2)}}{n^2}=\frac74\zeta(4),\qquad
 \sum_{n\ge1}\frac{H_n^2}{n^2}=\frac{17}4\zeta(4).
 \label{bhr:quartic-harmonic:euler}
\end{equation}
These are classical ordinary convergent sums. One short derivation of
the three weight-four formulas is as follows. The shuffle and stuffle
products of $\zeta(2)^2$ give respectively
$2\zeta(2,2)+4\zeta(3,1)$ and $2\zeta(2,2)+\zeta(4)$;
hence $\zeta(3,1)=\zeta(4)/4$ and
$\zeta(2,2)=3\zeta(4)/4$.
The iterated-integral substitution $t\mapsto1-t$ gives
$\zeta(2,1,1)=\zeta(4)$. Expanding the three harmonic numerators
in strict sums now gives exactly the displayed values. The first,
weight-three formula follows from the same shuffle/stuffle method or
from the elementary symmetric integral
$\int_0^1\log x\log(1-x)/(x(1-x))\,dx=2\zeta(3)$.

The decisive finite difference is
\begin{equation}
 r_n-r_{n-1}=\frac{12a_{n-1}}n+\frac4{n^3}.
 \label{bhr:quartic-harmonic:residue-difference}
\end{equation}
It follows by substituting
$b_n=b_{n-1}+n^{-1}$, $c_n=c_{n-1}+n^{-2}$,
and $e_n=e_{n-1}+n^{-3}$; no limiting argument is involved.
Thus finite summation by parts gives
\begin{align}
 \sum_{n=1}^N r_n\log\left(1+\frac1n\right)
 &=r_N\log(N+1)-12\sum_{n=1}^N\frac{a_{n-1}\log n}{n}
                      -4\sum_{n=1}^N\frac{\log n}{n^3}.
 \label{bhr:quartic-harmonic:SBP}
\end{align}

We also need the constant term in $\sum_{n=1}^N r_n/n$.
Here is an exact finite identity that makes this normalization
checkable. Put $X=H_N$, $Y=H_N^{(2)}$, $Z=H_N^{(3)}$,
$W=H_N^{(4)}$, and
\[
 S_{12}=\sum_{n=1}^N\frac{H_n}{n^2},\quad
 S_{13}=\sum_{n=1}^N\frac{H_n}{n^3},\quad
 S_{22}=\sum_{n=1}^N\frac{H_n^{(2)}}{n^2},\quad
 S_{112}=\sum_{n=1}^N\frac{H_n^2}{n^2}.
\]
Then
\begin{align}
 \sum_{n=1}^N\frac{r_n}{n}
 ={}&-X^4+4\gamma X^3+(-6\gamma^2+6\zeta(2))X^2
       +(4\gamma^3-12\gamma\zeta(2)+4\zeta(3))X\notag\\
 &+6X^2Y-12\gamma XY-4XZ
       +(-6\gamma^2+6\zeta(2))Y-16\gamma Z-11W\notag\\
 &+24\gamma S_{12}-12S_{112}+20S_{13}+6S_{22}.
 \label{bhr:quartic-harmonic:finite-sum}
\end{align}
For example this identity follows by taking finite differences of
$H_n^4$, $H_n^3$, $H_n^2H_n^{(2)}$, $H_nH_n^{(2)}$,
and $H_nH_n^{(3)}$, then collecting terms. Both sides vanish at
$N=0$, and their exact finite differences are $r_N/N$.

Let $L=\log N$. Using \eqref{bhr:quartic-harmonic:euler} in
\eqref{bhr:quartic-harmonic:finite-sum}, together with
$H_N=L+\gamma+O(N^{-1})$, gives
\begin{equation}
 \sum_{n=1}^N\frac{r_n}{n}
 =-L^4+12\zeta(2)L^2+C+o(1),\qquad
 C=\gamma^4-18\gamma^2\zeta(2)+32\gamma\zeta(3)
       -\frac{23}2\zeta(4).
 \label{bhr:quartic-harmonic:C}
\end{equation}
Also
$r_N\log(N+1)=-4L^4+24\zeta(2)L^2+o(1)$.
Since $a_{n-1}=2\zeta(2)-\log^2n+O(\log n/n)$, the limit
\begin{equation}
 A=\lim_{N\to\infty}\left\{
  \sum_{n=1}^N\frac{a_{n-1}\log n}{n}
      -\zeta(2)L^2+\frac{L^4}{4}\right\}
 \label{bhr:quartic-harmonic:A}
\end{equation}
exists. Equations \eqref{bhr:quartic-harmonic:K-eq},
\eqref{bhr:quartic-harmonic:SBP}, and \eqref{bhr:quartic-harmonic:C}
now imply
\begin{equation}
 Q_4=18\zeta(4)+12\zeta(3)-12\gamma\zeta(2)
                         +4\zeta'(3)-12A.
 \label{bhr:quartic-harmonic:before-Laurent}
\end{equation}
Every divergent polynomial has been retained until this cancellation.

\subsection{Convergent arithmetic and the Laurent coefficient}

At a positive integer,
\[
 a_{n-1}=2\zeta(2)-\psi(n)^2-\psi'(n).
\]
Moreover $\psi(n)^2+\psi'(n)-\log^2n=O(\log n/n)$.
Using the standard Stieltjes defining limits in
\eqref{bhr:quartic-harmonic:A} gives
\[
 A=2\zeta(2)\gamma_1-\gamma_3
   -\sum_{n=1}^{\infty}
          \frac{\bigl(\psi(n)^2+\psi'(n)-\log^2n\bigr)\log n}{n}.
\]
The final series is absolutely convergent because its summand is
$O(\log^2n/n^2)$. Substitution proves
\eqref{bhr:quartic-harmonic:arithmetic}.

To identify the harmonic coordinate, introduce
\[
 F(s)=\sum_{n=1}^{\infty}
    \frac{\psi(n)^2+\psi'(n)-\log^2n}{n^s},\qquad\Re s>0.
\]
This series and every fixed spectral derivative converge normally
on compact subsets of the stated half-plane. Expanding the harmonic
numerator in strict sums gives, initially for $\Re s>1$,
\begin{equation}
 F(s)=2Z_3(s,1,1)-2\gamma Z_2(s,1)
                   +(\gamma^2+\zeta(2))\zeta(s)-\zeta''(s).
 \label{bhr:quartic-harmonic:deformation}
\end{equation}
This also supplies the required local continuation. The coefficient
of $u$ at $s=1+u$ is
\[
 F'(1)=2\beta_1-2\gamma\eta
                 -(\gamma^2+\zeta(2))\gamma_1+\gamma_3.
\]
On the other hand it is the negative of the convergent sum in
\eqref{bhr:quartic-harmonic:arithmetic}. Inserting this relation proves
\eqref{bhr:quartic-harmonic:main-eq}. Combining it with the cubic bridge
proves \eqref{bhr:quartic-harmonic:combined}.

\subsection{A Mellin normalization for every height-one coefficient}

The following classical Mellin mechanism (compare the harmonic-zeta framework in \cite{bhr:Kargin}) supplies an independent
realization of the coefficient used above. Its all-depth formulation
also states precisely where new global information first enters.
Let $d\ge1$, and write
\[
 Z_d(s,\{1\}^{d-1})
   =\sum_{n=1}^{\infty}\frac{H_{n-1}(\{1\}^{d-1})}{n^s},
 \qquad
 \frac1{\Gamma(1+u)}=\sum_{j\ge0}c_ju^j.
\]
An empty harmonic word has value one. For $t>0$ put
\[
 B(t)=-\log(1-e^{-t}),\qquad
 V_d(t)=B(t)^d-\mathbf1_{0<t<1}(-\log t)^d.
\]
Since $V_d(t)=O(t(1+|\log t|^{d-1}))$ at zero and decays
exponentially at infinity, every logarithmic moment of $V_d(t)/t$
converges absolutely.

\begin{lemma}[All height-one finite parts and regular coefficients]
\label{bhr:quartic-harmonic:mellin}
For $|u|<1$, as an identity of meromorphic functions,
\begin{equation}
 \Gamma(1+u)Z_d(1+u,\{1\}^{d-1})
   =\frac1{u^d}+\frac{u}{d!}
       \int_0^\infty t^{u-1}V_d(t)\,dt.
 \label{bhr:quartic-harmonic:mellin-generator}
\end{equation}
Consequently the principal part and constant coefficient together are
$\sum_{j=0}^{d}c_j u^{j-d}$. For every $k\ge1$, the exact regular
coefficient is
\begin{equation}
 [u^k]Z_d(1+u,\{1\}^{d-1})
  =c_{d+k}+\frac1{d!}\sum_{j=1}^k
       \frac{c_{k-j}}{(j-1)!}
       \int_0^\infty\frac{V_d(t)\log^{j-1}t}{t}\,dt.
 \label{bhr:quartic-harmonic:mellin-coefficients}
\end{equation}
In particular the finite part is $c_d$, while the coefficient of $u$
is $c_{d+1}+d!^{-1}\int_0^\infty V_d(t)\,dt/t$.
\end{lemma}

\begin{proof}
The elementary multiple-polylogarithm identity
$\operatorname{Li}_{\{1\}^{d-1}}(z)
=(-\log(1-z))^{d-1}/(d-1)!$ gives, first for $\Re s>1$,
\[
 \Gamma(s)Z_d(s,\{1\}^{d-1})
  =\int_0^\infty t^{s-1}
       \frac{B(t)^{d-1}}{(d-1)!(e^t-1)}\,dt.
\]
Subtract $(-\log t)^{d-1}/((d-1)!t)$ on $(0,1)$.
Its Mellin integral at $s=1+u$ is $u^{-d}$. The subtracted
kernel is $-V_d'(t)/d!$; no jump term occurs at $t=1$, since
$(-\log1)^d=0$. Integrating that remainder by parts produces
exactly the integral in \eqref{bhr:quartic-harmonic:mellin-generator}.
The boundary terms vanish for $\Re u>-1$, because of the bounds on
$V_d$. Normal convergence of that integral and all its spectral
derivatives proves continuation and coefficient extraction.
\end{proof}

For the present depth-three coordinate this yields the particularly
simple convergent formula
\begin{equation}
 \beta_1=
 \frac{\gamma^4}{24}-\frac{\gamma^2\zeta(2)}4
       +\frac{\gamma\zeta(3)}3+\frac{\zeta(4)}{16}
       +\frac16\int_0^\infty\frac{V_3(t)}t\,dt.
 \label{bhr:quartic-harmonic:beta-mellin}
\end{equation}
Likewise $\eta=c_3+\tfrac12\int_0^\infty V_2(t)\,dt/t$.
These formulas evaluate the harmonic coefficients through elementary
convergent kernels; they do not identify those kernels with finite
combinations of ordinary constants. They also confirm the constant
coefficient displayed in \eqref{bhr:quartic-harmonic:beta} directly.

\subsection{A normalized primitive and every differentiated identity}

Put
\[
 d_0=6\gamma^2-4\zeta(2),\qquad
 r_0=4\gamma^3-12\gamma\zeta(2)+4\zeta(3).
\]
For $\Re z>0$ with the principal logarithm, define
\begin{align}
 \mathcal P_4(z)={}&-\frac1{3z^3}-\frac{2\gamma}{z^2}
                      -\frac{d_0}{z}+r_0\log z+C_*z\notag\\
 &+\sum_{n=1}^{\infty}\Bigg[
   \frac1{3n^3}-\frac1{3(n+z)^3}-\frac z{n^4}\notag\\
 &\hspace{11mm}-4b_n\left\{\frac1{2n^2}-\frac1{2(n+z)^2}
                                      -\frac z{n^3}\right\}\notag\\
 &\hspace{11mm}+d_n\left\{\frac1n-\frac1{n+z}-\frac z{n^2}\right\}
       +r_n\left\{\log\left(1+\frac zn\right)-\frac zn\right\}
                                      \Bigg].
 \label{bhr:quartic-harmonic:primitive}
\end{align}
The normally convergent braces are the integrals, vanishing at zero,
of the subtracted principal parts. Consequently
\[
 \mathcal P_4'(z)=\psi(z)^4,\qquad
 \operatorname{FP}\int_0^b\psi(x)^4\,dx=\mathcal P_4(b)
      \quad(b>0),\qquad \mathcal P_4(1)=Q_4.
\]
The constant term of $\mathcal P_4$ at the singular origin is zero;
this fixes the integration constant without an unspecified branch term.

For every integer $m\ge1$, differentiation of
\eqref{bhr:quartic-harmonic:ML-eq} gives the further exact identity
\begin{align}
 &\sum_{m_1+m_2+m_3+m_4=m}
   \frac{m!}{m_1!m_2!m_3!m_4!}
      \prod_{j=1}^4\psi^{(m_j)}(z)\notag\\
 &\quad=(-1)^m m!\sum_{n=0}^{\infty}
 \left\{
  \frac{\binom{m+3}{3}}{(n+z)^{m+4}}
  -\frac{4b_n\binom{m+2}{2}}{(n+z)^{m+3}}
  +\frac{(m+1)d_n}{(n+z)^{m+2}}
  +\frac{r_n}{(n+z)^{m+1}}
 \right\}.
 \label{bhr:quartic-harmonic:derivatives}
\end{align}
The convergence is normal off the digamma poles, since the slowest
term is $O(\log^3n/n^{m+1})$. At $z=1$, the left side is a finite
polynomial in $\gamma$ and ordinary zeta values.

\paragraph{Independent checks.}
The exact checking program verifies the principal parts, the residue
finite difference, the full finite identity
\eqref{bhr:quartic-harmonic:finite-sum}, and the cancellations producing
\eqref{bhr:quartic-harmonic:main-eq}. A separate numerical calculation
integrates a locally Taylor-subtracted digamma fourth power and compares
it with the convergent arithmetic formula, whose tail is evaluated by
Bernoulli asymptotics and Hurwitz-zeta derivatives. At working precision
80 digits, the arithmetic-tail evaluation differs from quadrature by
less than $3.4\cdot10^{-62}$. Independent Mellin evaluations of $\eta$
and $\beta_1$, inserted into the main identity, differ from the same
quadrature by less than $2.1\cdot10^{-74}$. The common value begins
\[
 Q_4=13.5101203760158274116062407128304484419327985333019186009860964\ldots.
\]
These are numerical diagnostics, not certified enclosures or substitutes
for the proof.

\paragraph{Next precise targets.}
The theorem determines $\beta_1$ from the cubic and quartic finite parts,
but does not express it in a prescribed algebra of ordinary constants.
For the fifth power, the residue polynomial has degree four in the
finite harmonic data, and its first difference should produce both
height-one depth-four Laurent data and additional harmonic numerators.
The correct next task is to derive that exact difference and identify the
minimal \emph{stated} family of Laurent coefficients. A second direction
is a common Hurwitz shift: the shifted strict harmonic numerator and the
integration endpoint must be normalized together, since replacing
$n$ by $n+a$ only in the outer denominator does not shift the harmonic
prefix or the finite-part coordinate.

% END sections/05_quartic_harmonic.tex


% BEGIN sections/06_cyclic_transverse.tex
% Integration fragment: requires amsmath, amssymb, amsthm and theorem/proposition environments.
% Labels use fcy:; bibliography keys are shared with the full article.
\section{An ordinary convergent generator for symmetric Tornheim jets}
\label{bhr:fcy:section}

The incoming \emph{Harmonic Parity and Resolvent Identities} report
already determines the quartic directional layer and the formal
all-order kernel. Its Question R10 asks for additional
analytic information about the symmetric quadratic germ controlling
cyclic fifth derivatives. We answer the stated integral version of
that question. We first give an explicit, normally convergent integral
for its full symmetric generating germ. We then extract one
transverse quadratic coefficient, including its full finite
normalization. Together with the diagonal fifth derivative this
coefficient determines every cyclic fifth derivative.

No claim that either retained coefficient has a finite evaluation in
ordinary zeta constants is made. The earlier quartic reconstruction,
formal dimension count, and cyclic factorization are used with their
existing attribution \cite{bhr:Parity}. The holomorphy of the cyclic
Tornheim sum near the origin is also a special case of Nakamura's
symmetric desingularization theorem \cite[Theorem 1.2]{bhr:Nakamura}.

\subsection{The fixed holomorphic remainder}

Write
\[
 T(A,B,C)=\sum_{m,n\ge1}\frac1{m^A n^B(m+n)^C},\qquad
 \mathfrak g(C)=\frac1{\Gamma(1+C)},\qquad \Sigma=A+B+C.
\]
The series is first used where it converges absolutely; subsequently
\(T\) denotes its meromorphic continuation. We use exactly the
normalization of the preceding report:
\begin{equation}\label{bhr:fcy:R}
 T(A,B,C)=\zeta(A)\zeta(B)
 -\frac{C\zeta(B-1)}{A+C}-\frac{C\zeta(A-1)}{B+C}
 +C R(A,B,C),
\end{equation}
where \(R\) is holomorphic and symmetric in \(A,B\) near the
origin. Here and below an apparent quotient of holomorphic germs by a
linear form is assigned its holomorphic extension when its numerator
is divisible by that form.

The finite elementary kernel needed below is
\begin{align}
 E(A,B,C)={}&\mathfrak g(C)\left\{
 \frac{\Gamma(1-A)\Gamma(1-B)}{\Sigma-2}
 +\frac{\Gamma(1-A)\zeta(B)}{A+C-1}
 +\frac{\Gamma(1-B)\zeta(A)}{B+C-1}\right\}\notag\\
 &+\zeta(A)\zeta(B)\frac{\mathfrak g(C)-1}{C}\notag\\
 &-\zeta(B-1)\frac{\Gamma(1-A)\mathfrak g(C)-1}{A+C}
 -\zeta(A-1)\frac{\Gamma(1-B)\mathfrak g(C)-1}{B+C}.
 \label{bhr:fcy:E}
\end{align}
It is jointly holomorphic near zero. For example,
\(\Gamma(1-A)\mathfrak g(-A)=1\), proving removability at \(A+C=0\).

For completeness, put \(F_A(x)=\operatorname{Li}_A(e^{-x})\) and
\begin{align}
 S(A,B;x)={}&\Gamma(1-A)\Gamma(1-B)x^{A+B-2}
 +\Gamma(1-A)\zeta(B)x^{A-1}
 +\Gamma(1-B)\zeta(A)x^{B-1}\notag\\
 &-\Gamma(1-A)\zeta(B-1)x^A
 -\Gamma(1-B)\zeta(A-1)x^B+\zeta(A)\zeta(B).
 \label{bhr:fcy:S}
\end{align}
Jonqui\`ere's expansion \cite[25.12.12]{bhr:DLMFPolylog} gives the
normally convergent integral
\begin{align}
 I_0(A,B,C)={}&\int_0^1x^{C-1}\bigl(F_A(x)F_B(x)-S(A,B;x)\bigr)\,dx
 \notag\\
 &+\int_1^\infty x^{C-1}F_A(x)F_B(x)\,dx
 \label{bhr:fcy:I0}
\end{align}
and the identity
\begin{equation}\label{bhr:fcy:RE}
 R(A,B,C)=\mathfrak g(C)I_0(A,B,C)+E(A,B,C).
\end{equation}
Indeed, integrate the six subtracted terms explicitly in the Mellin
representation for \(T\), multiply by \(\mathfrak g(C)\), and collect the
two polar terms of \eqref{bhr:fcy:R}. This gives
\eqref{bhr:fcy:E} term by term. In a sufficiently small closed parameter
polydisc the omitted Jonqui\`ere terms, after multiplication by
\(x^{C-1}\), are bounded by
\(x^{-\delta}(1+|\log x|^N)\) for some \(\delta<1\), including
any prescribed finite number of parameter derivatives. At infinity
there is uniform exponential decay. These estimates justify the
joint holomorphy and all differentiations used below.

\subsection{An explicit integral for the symmetric germ}

The earlier factorization defines the unique symmetric holomorphic
germ \(J\) by
\begin{align}
 \sum_{\mathrm{cyc}}T(A,B,C)={}&
 2\bigl(\zeta(A)\zeta(B)+\zeta(B)\zeta(C)+\zeta(C)\zeta(A)\bigr)
 \notag\\
 &-\zeta(A+B)-\zeta(B+C)-\zeta(C+A)\notag\\
 &+2\bigl(\zeta(A)+\zeta(B)+\zeta(C)\bigr)+1+ABCJ(A,B,C).
 \label{bhr:fcy:J}
\end{align}
We retain this normalization exactly. The next formula expresses
\(J\) through ordinary polylogarithms, Gamma and zeta functions,
and ordinarily convergent integrals, without any Tornheim derivative
on the right-hand side.

For a holomorphic function \(H\), define
\[
 \mathcal D_{A,B}H(A,B,C)=
 \frac{H(A,B,C)-H(A,0,C)-H(0,B,C)+H(0,0,C)}{AB}.
\]
At zero parameters this means the corresponding derivative. Set
\begin{align*}
 d_A(x)&=\frac{\operatorname{Li}_A(e^{-x})-\operatorname{Li}_0(e^{-x})}{A},&
 g_A(x)&=\frac{\Gamma(1-A)x^A-1}{A},\\
 p_A&=\frac{\zeta(A)+1/2}{A},&
 q_A&=\frac{\zeta(A-1)+1/12}{A},
\end{align*}
with the holomorphic values at \(A=0\). Define
\begin{align}
 S_\Delta(A,B;x)={}&\frac{g_A(x)g_B(x)}{x^2}
 +\frac{g_A(x)p_B+g_B(x)p_A}{x}\notag\\
 &-g_A(x)q_B-g_B(x)q_A+p_Ap_B,
 \label{bhr:fcy:SDelta}\\
 I_\Delta(A,B,C)={}&\int_0^1x^{C-1}
 \bigl(d_A(x)d_B(x)-S_\Delta(A,B;x)\bigr)\,dx\notag\\
 &+\int_1^\infty x^{C-1}d_A(x)d_B(x)\,dx.
 \label{bhr:fcy:IDelta}
\end{align}

\begin{theorem}[Ordinary convergent symmetric generator]
\label{bhr:fcy:generator}
The integrals \eqref{bhr:fcy:IDelta} converge normally in a common
neighborhood of \((0,0,0)\), including after arbitrary fixed
parameter derivatives, and
\begin{equation}\label{bhr:fcy:Jintegral}
 J(A,B,C)=\sum_{\mathrm{cyc}}
 \left\{\mathfrak g(C)I_\Delta(A,B,C)+\mathcal D_{A,B}E(A,B,C)\right\}.
\end{equation}
\end{theorem}

\begin{proof}
Apply the rectangular difference in \(A,B\) to
\eqref{bhr:fcy:I0}. The difference of the product \(F_AF_B\),
divided by \(AB\), is \(d_Ad_B\); the corresponding difference
of \(S\) is exactly \(S_\Delta\). Cauchy estimates on a slightly
larger parameter polydisc preserve the integrable bounds of
\eqref{bhr:fcy:I0}; equivalently, one can use
\(d_A(x)=g_A(x)/x+p_A-q_Ax+O(x^{2-\epsilon})\), locally
uniformly with logarithmic factors after differentiation. Thus
\(I_\Delta=\mathcal D_{A,B}I_0\), with normal convergence.

In the cyclic sum of \eqref{bhr:fcy:R}, the two terms with
denominator \(A+C\) combine to \(-\zeta(B-1)\), and similarly
for the other denominators. All remaining explicit zeta terms
depend on at most two variables. The complete rectangular
three-variable difference therefore annihilates them and gives
\[
 ABCJ(A,B,C)=\sum_{\mathrm{cyc}}C\bigl(
 R(A,B,C)-R(A,0,C)-R(0,B,C)+R(0,0,C)\bigr).
\]
To check the left-hand side, apply the same difference to
\eqref{bhr:fcy:J}: all its explicit terms vanish, and \(ABCJ\)
vanishes on every coordinate plane. Divide by \(ABC\), substitute
\eqref{bhr:fcy:RE}, and use holomorphic extension at zero parameters.
This proves \eqref{bhr:fcy:Jintegral}.
\end{proof}

The theorem is an analytic representation of the previously defined
germ. It supplies all symmetric jet coefficients in one fixed,
ordinary convergent kernel. Its content includes the endpoint
subtractions and their finite normalization; replacing the kernel by
an unsubtracted integral would not preserve its Taylor coefficients.

\subsection{The first transverse quadratic coefficient}

Write \(\gamma\) for Euler's constant and put
\[
 L=\log(2\pi),\qquad z_j=\zeta^{(j)}(0),\qquad
 y_j=\zeta^{(j)}(-1),\qquad \kappa_2=\zeta(2),\quad \kappa_4=\zeta(4).
\]
The quadratic homogeneous part of \(J\) has the form
\begin{equation}\label{bhr:fcy:J2form}
 J_2(A,B,C)=j_{20}(A^2+B^2+C^2)+\chi(AB+BC+CA).
\end{equation}
Thus \(\chi=[AB]J\) is the additional scalar to be determined.
For \(x>0\), let
\[
 f_j(x)=\left.\partial_s^j\operatorname{Li}_s(e^{-x})\right|_{s=0},
 \qquad h(x)=\log x+\gamma,
\]
and define
\begin{align}
 P_{21}(x)={}&\frac{h(h^2+\kappa_2)}{x^2}
 +\frac{z_1(h^2+\kappa_2)+z_2h}{x}
 -y_1(h^2+\kappa_2)-y_2h+z_1z_2,
 \label{bhr:fcy:P21}\\
 P_{22}(x)={}&\frac{(h^2+\kappa_2)^2}{x^2}
 +\frac{2z_2(h^2+\kappa_2)}{x}-2y_2(h^2+\kappa_2)+z_2^2,
 \label{bhr:fcy:P22}\\
 \mathcal K(x)={}&\frac14 f_2(x)^2+h(x)f_2(x)f_1(x),\qquad
 P(x)=\frac14P_{22}(x)+h(x)P_{21}(x).
 \label{bhr:fcy:KP}
\end{align}
Let
\begin{equation}\label{bhr:fcy:chiintegral}
 \mathcal I=\int_0^1\frac{\mathcal K(x)-P(x)}x\,dx
                   +\int_1^\infty\frac{\mathcal K(x)}x\,dx.
\end{equation}
Both integrals converge ordinarily and absolutely.

\begin{theorem}[Explicit transverse fifth-order coordinate]
\label{bhr:fcy:chi}
The coefficient in \eqref{bhr:fcy:J2form} is
\begin{equation}\label{bhr:fcy:chiresult}
 \chi=\mathcal I+\mathcal E,
\end{equation}
where
\begin{align}
 \mathcal E={}&-\frac{10\gamma^4+20\gamma^3+12\gamma^2\kappa_2
 +30\gamma^2+12\gamma \kappa_2+30\gamma+2\kappa_2^2+6\kappa_2+15}{16}
 \notag\\
 &+\frac{-\gamma^4-2\gamma^2\kappa_2+\kappa_2^2+2\kappa_4}{4}\,y_1
 -\frac{\gamma(\gamma^2+\kappa_2)}2\,y_2\notag\\
 &-\bigl(\gamma^3+3\gamma^2+6\gamma+6+(\gamma+1)\kappa_2\bigr)z_1
 -\frac{3\gamma^2+6\gamma+6+\kappa_2}{2}\,z_2\notag\\
 &+\frac{\gamma^2-\kappa_2}{2}\,z_1z_2+\frac\gamma4z_2^2.
 \label{bhr:fcy:chicorrection}
\end{align}
\end{theorem}

\begin{proof}
All derivatives of \(R\) in this proof are evaluated at the
origin. The coefficient of \(A^2B^2C\) in
\(\sum_{\mathrm{cyc}}CR(A,B,C)\) is
\[
 \frac14R_{AABB}+R_{AABC}.
\]
The first cyclic term contributes \(R_{AABB}/4\), and the
other two contribute \(R_{AABC}/2\) each. The explicit
zeta terms in \eqref{bhr:fcy:J} and the cyclic form of
\eqref{bhr:fcy:R} have no monomial involving all three variables.
Consequently
\begin{equation}\label{bhr:fcy:chiderivative}
 \chi=\frac14R_{AABB}+R_{AABC}.
\end{equation}

Differentiate \eqref{bhr:fcy:RE}, using \(\mathfrak g(0)=1\) and
\(\mathfrak g'(0)=\gamma\). The integral contribution to
\eqref{bhr:fcy:chiderivative} is precisely \eqref{bhr:fcy:chiintegral}:
\(S_{AAB}=P_{21}\), \(S_{AABB}=P_{22}\), and the derivative
of \(\mathfrak g(C)x^{C-1}\) at zero is \(h(x)/x\).
For a direct convergence check, the differentiated Jonqui\`ere
expansions are
\[
 f_1(x)=\frac h x+\sum_{k\ge0}
       \frac{(-1)^k\zeta'(-k)}{k!}x^k,\qquad
 f_2(x)=\frac{h^2+\kappa_2}x+\sum_{k\ge0}
       \frac{(-1)^k\zeta''(-k)}{k!}x^k.
\]
The displayed subtractions remove both singular products, the
cross terms of analytic degrees zero and one, and the constant
analytic product. Hence \((\mathcal K-P)/x=O(1+|\log x|^3)\) at zero.
At infinity \(f_1,f_2=O(e^{-2x})\), so
\(\mathcal K/x=O((1+\log x)e^{-4x}/x)\).

The remaining finite contribution is
\[
 \mathcal E=[A^2B^2]E+2[A^2BC]E.
\]
Insert
\begin{align*}
 \log\Gamma(1-A)&=\gamma A+\frac{\kappa_2}{2} A^2
                  +\frac{\zeta(3)}3A^3+\frac{\kappa_4}4A^4+O(A^5),\\
 \log \mathfrak g(C)&=\gamma C-\frac{\kappa_2}{2} C^2
                  +\frac{\zeta(3)}3C^3-\frac{\kappa_4}4C^4+O(C^5)
\end{align*}
into the elementary expression \eqref{bhr:fcy:E}. Expand its three
nonresonant denominators geometrically and divide the three
removable quotients as holomorphic power series. Extraction of
the two coefficients gives \eqref{bhr:fcy:chicorrection}. This is
finite polynomial algebra; the accompanying exact script performs
it independently with all zeta derivatives left symbolic.
\end{proof}

\subsection{Every cyclic fifth derivative}

Put \(\Omega_5=W_{1,1,1}^{(5)}(0)\), where
\(W_{a,b,c}(t)=T(at,bt,ct)\) is restricted before taking its
one-variable continuation at zero. Write
\[
 e_1=a+b+c,\quad e_2=ab+bc+ca,\quad e_3=abc,\quad
 D=e_1^2-3e_2,\quad p_2=e_1^2-2e_2.
\]

\begin{theorem}[Cyclic fifth-order formula]\label{bhr:fcy:fifth}
If the three pairwise sums of \(a,b,c\) are nonzero, then
\begin{align}
 \sum_{\mathrm{cyc}}W_{a,b,c}^{(5)}(0)={}&
 e_3p_2\Omega_5-120e_3D\chi
 -5(e_1e_2-3e_3)D\,Lz_4\notag\\
 &+20(e_1e_2^2-4e_1^2e_3+3e_2e_3)z_2z_3\notag\\
 &+(-2e_1^5+5e_1^3e_2-5e_1e_2^2
                         +47e_1^2e_3-69e_2e_3)z_5.
 \label{bhr:fcy:fifthformula}
\end{align}
For arbitrary complex slopes the same polynomial gives the fifth
derivative of the holomorphic cyclic sum restricted as a whole.
This last assertion does not assign values to its individual
singular summands.
\end{theorem}

\begin{proof}
On the diagonal, differentiation of \eqref{bhr:fcy:J} gives
\[
 j_{20}+\chi=\frac{
 \Omega_5+10Lz_4-40z_2z_3+32z_5}{120}.
\]
Thus \eqref{bhr:fcy:J2form} is completely determined by
\(\Omega_5\) and the ordinarily convergent coordinate
\eqref{bhr:fcy:chiresult}. For general slopes, the fifth derivative
of the explicit part of \eqref{bhr:fcy:J} is
\[
 -5Lz_4\sum_{\mathrm{sym}}a^4b
 +20z_2z_3\sum_{\mathrm{sym}}a^3b^2
 -\bigl((a+b)^5+(b+c)^5+(c+a)^5\bigr)z_5.
\]
The remaining derivative is \(120abcJ_2(a,b,c)\).
Substitute the preceding value of \(j_{20}\), use the elementary
symmetric polynomial identities, and collect terms. This yields
\eqref{bhr:fcy:fifthformula}. Since \eqref{bhr:fcy:J} is jointly
holomorphic, its cyclic restriction exists even when an individual
pairwise divisor would prevent restricting a single Tornheim term.
\end{proof}

For instance, the explicit asymmetric identity is
\begin{align}
 &W_{1,2,3}^{(5)}(0)+W_{2,3,1}^{(5)}(0)+W_{3,1,2}^{(5)}(0)
       -84\Omega_5\notag\\
 &\hspace{18mm}=-2160\chi-720L\zeta^{(4)}(0)
       +1200\zeta''(0)\zeta'''(0)-1704\zeta^{(5)}(0).
 \label{bhr:fcy:example}
\end{align}
For real slopes, \(D=\tfrac12((a-b)^2+(b-c)^2+(c-a)^2)\).
Consequently a single real cyclic fifth ray with \(abc\ne0\)
eliminates the displayed transverse coefficient only on the
diagonal. For complex slopes the cone \(D=0\) has additional
points. These are coefficient statements in the specified two
coordinate representation, not claims of arithmetic independence.

\needspace{13\baselineskip}
\subsection{Numerical checks and the remaining arithmetic problem}
The package evaluates \(\mathcal I\) using differentiated
Jonqui\`ere series on \((0,1)\) and exponentially convergent
incomplete-Gamma sums on \((1,\infty)\). At 70 decimal digits
of working precision the diagnostic values are
\begin{align*}
 \mathcal I&=-0.0128358838132664551878307355997471413775071472849\ldots,\\
 \mathcal E&=18.6522612153652911636439613245591942979227651526949\ldots,\\
 \chi&=18.6394253315520247084561305889594471565452580054100\ldots,\\
 \Omega_5&=8808.87952093843854353140320772504731742914386526227\ldots.
\end{align*}
An independently evaluated Mellin continuation of individual
\(T\) rays checks \eqref{bhr:fcy:fifthformula} on the cycles generated
by \((1,2,3)\), \((1,1,2)\), and \((1,-2,3)\); the largest
recorded discrepancy is below \(6\cdot10^{-30}\). These are
uncertified numerical diagnostics, not interval enclosures or
proofs of the analytic identities. The exact coefficient script
checks the elementary correction and the full symbolic fifth-order
formula.

The new generator answers the ordinary-integral alternative of
Question R10 and makes all higher symmetric coefficients concrete.
Finite ordinary-constant evaluations of \(\chi\) and
\(\Omega_5\) remain open. A sharper next problem is to find an
independent Gamma or Barnes representation of the kernel in
\eqref{bhr:fcy:chiintegral}, with its normalization \eqref{bhr:fcy:chicorrection}
unchanged. At the following order, the same generator yields the
three coefficients of the symmetric cubic germ \(J_3\); any
reduction below that formal count requires an additional identity.

% END sections/06_cyclic_transverse.tex


% BEGIN sections/07_audit_research.tex
\section{Source audit, verification, and further research}
\label{bhr:sec:audit}

\subsection{A precise correction to a published Barnes identity}

The inspection of the Barnes Fourier source revealed a pointwise
algebraic error in equation (7) of the inspected arXiv version of
Mez\H{o}'s paper \cite{bhr:Mezo}. That equation identifies
$\log^2G(x)$ with
\[
 \frac14\left\{
   \log^2\bigl(G(x)G(1-x)\bigr)
   +\log^2\bigl(G(x)/G(1-x)\bigr)\right\}.
\]
On $(0,1)$ both Barnes values are positive. Put
$a=\log G(x)$ and $b=\log G(1-x)$. The displayed expression
is $(a^2+b^2)/2$, not $a^2$. Thus the corrected pointwise identity is
\begin{equation}\label{bhr:eq:Mezo-correction}
 \boxed{\displaystyle
 \frac{\log^2G(x)+\log^2G(1-x)}2
 =\frac14\left\{
   \log^2\bigl(G(x)G(1-x)\bigr)
   +\log^2\bigl(G(x)/G(1-x)\bigr)\right\}.}
\end{equation}
Alternatively, to keep the original left side one must add
\[
 \frac12\log\bigl(G(x)G(1-x)\bigr)
           \log\bigl(G(x)/G(1-x)\bigr)
\]
to the original right side.

This is an exact objection, not a small numerical discrepancy.
From $G(1+x)=\Gamma(x)G(x)$ and $G(1)=1$ one gets
$G(x)=x(1+O(x))$ as $x\downarrow0$, while
$\log G(1-x)=O(x)$. Therefore the original left side has leading
term $\log^2x$, whereas its stated right side has leading term
$\tfrac12\log^2x$.

The distinction has a narrow consequence. After integration over
$(0,1)$, reflection makes
$\int\log^2G(1-x)=\int\log^2G(x)$, and the added product
is odd under $x\mapsto1-x$. The integrated equality used next in
that paper is valid. We therefore propose correcting the pointwise
line or placing its equality after integration; we do not infer
that its subsequent integrated theorem is false. The source version
and location are recorded in \src{integration/CORRECTIONS.md}.

\subsection{What the repository audit does and does not establish}

The source inventory contains 81 manuscript chapter files and all
11 incoming archives present at the pinned commit. The detailed
mathematical review concentrates on the conjecture status records,
the cotangent-resolvent programme, the harmonic power bridges, and
the cyclic Tornheim remainder. The inventory is not a claim that
every theorem in all these documents has received a fresh proof audit.

The following distinctions are necessary for correct integration.
\begin{enumerate}[leftmargin=*,itemsep=5pt]
 \item The classical Barnes/Hurwitz conversion, the log-$G$ Fourier
 coefficients, balanced negapolygamma primitives, beta-integral
 binomial identities, and symmetric Tornheim desingularization keep
 their existing attribution. The new contribution is the proved
 resolvent and coefficient reduction with the stated normalization.
 \item The earlier report already contains the quartic directional
 Tornheim reconstruction and the formal cyclic dimension count.
 Section~\ref{bhr:fcy:section} supplies the requested explicit convergent
 integral, and then the fifth-order formula using that analytic
 coordinate; it does not relabel the earlier formal result as new.
 \item The quartic digamma moment in
 Section~\ref{bhr:quartic-harmonic:sec} is a different question from
 the quartic directional Tornheim layer. Its depth-three coefficient
 $\beta_1$ is strict, as specified in
 \eqref{bhr:quartic-harmonic:beta}. An inclusive harmonic coefficient
 cannot be substituted without its diagonal correction.
 \item At a resonance, a Laurent constant in the spectral or kernel
 exponent is not automatically the finite part in the original
 integration coordinate. Equations~\eqref{bhr:eq:weighted-contact}
 and \eqref{bhr:cxp:eq:lambdaCT} explicitly perform the conversion.
 \item The current $S_6$ and revised $S_8$ coefficient vectors remain
 conjectural. The pre-existing rejection of the old frozen $S_8$
 vector must retain its separate label and status. No present
 decimal evaluation changes any of these statuses.
\end{enumerate}

Within the inspected repository material we found no additional
mathematical error that is established by the present analysis.
The pointwise Barnes correction above is to the cited external
source. The attached claim ledger distinguishes proved identities,
classical ingredients, numerical diagnostics, and open questions.
The package consists of research documents and reproducible
calculations, not machine-checked ProveIt proof objects.

\subsection{Independent checks and reproducibility}

Four independent lines of calculation accompany the proofs.
The Barnes checks evaluate the depth-two derivative from its own
double series and evaluate the Fourier sum $\sum q^nT_n$ separately.
The same code compares the complete transform with direct
$\log G$ quadrature in both half-planes, and checks the removable
$q=0$ real component. The polynomial Hurwitz checks compare the
subtracted double series and its spectral derivative with direct
Hurwitz quadrature, including a weighted Stieltjes coefficient.

The complex-power checks compare its Mellin generator with the
original integral at nonintegral complex parameters, and independently
test its Gamma closure and a noncommuting finite-part example.
The quartic checks use three routes: local digamma subtraction,
the convergent harmonic series with a zeta-derivative tail, and
the independent Mellin coefficients $\eta,\beta_1$.
The cyclic fifth checks compare the new ordinary-integral coordinate
with an independently evaluated continuation of three asymmetric
cycles. The pre-existing evaluator used in that comparison is
preserved with its provenance.

Finite exact algebra checks cover the multiple-Gamma coefficient
polynomials and normalization, quartic principal parts and their
finite summation identity, the elementary Tornheim correction,
and the symbolic fifth-order polynomial. These calculations are
useful for detecting signs, factorials, and omitted constants.
The analytic assertions---normal convergence, uniqueness of an
entire remainder, branches, continuation, and the permitted
coefficient interchanges---are proved in the text.

All numerical records are explicitly non-interval diagnostics.
Their working precision and truncations are recorded. Several
independent values agreeing to many digits cannot certify an
integer relation or a finite closed form. Scripts, dependency
versions, replay commands, numerical records, and a file-hash
manifest are included so that integration can preserve both the
theorem and its actual verification status.

\subsection{Further research questions}

The next problems are formulated as exact identities or reductions
in specified function classes. The first four build most directly
on the new formulas.

\begin{enumerate}[label=\textbf{F\arabic*.},leftmargin=*,itemsep=9pt]
 \item \textbf{Products of individual Barnes logarithms.}
 Determine the minimal explicit nested-series class for
 \[
  \int_0^1 W(x)\log\Gamma_r(x)\log\Gamma_t(x)R_z(x)\,dx.
 \]
 The polynomial Hurwitz decomposition reduces this to a bilinear
 weighted transform. The real question is whether its lattice
 splitting closes at depth three, with all diagonal corrections
 explicit, and which reflection combinations reduce the depth.

 \item \textbf{A second representation of $\mathcal D(q)$.}
 Find a Gamma, Barnes, or iterated-integral expression for
 \eqref{bhr:eq:Dq} that is independent of its defining depth-two sum.
 A useful target is an exact differential or functional equation
 under $q\mapsto q/(q-1)$, with a specified branch domain and
 initial value. Establish whether it reduces to ordinary
 polylogarithms and their order derivatives; do not infer such a
 reduction from a successful numerical fit at isolated points.

 \item \textbf{All higher digamma powers.}
 Starting from the principal parts of $\psi(z)^5$, derive the
 exact finite difference of its residue polynomial. Determine
 which height-one depth-four Laurent coefficients and which
 additional harmonic numerators survive summation by parts.
 A satisfactory result must state a finite family of retained
 coefficients and fix the entire remainder and every finite-part
 constant, as was done for the fourth power here.

 \item \textbf{Shifted quartic harmonic bridges.}
 For $a>0$, derive a common-shift analogue of
 \eqref{bhr:quartic-harmonic:main-eq}, using prefixes
 $\sum_{k=0}^{n-1}(k+a)^{-r}$ and the corresponding outer
 denominator. Track simultaneously the origin of the integration
 coordinate, the lower summation endpoint, and generalized
 Stieltjes constants $\gamma_j(a)$. Merely shifting the outer
 denominator leaves a different problem.

 \item \textbf{Nonpositive kernel resonances.}
 Extend the pointwise finite-part comparison of
 Theorem~\ref{bhr:cxp:thm:contacts} to $\Re\lambda\le0$,
 keeping contributions from both endpoints. There is already a
 concrete additional phase at $\lambda=-1$. The continued mean
 of the kernel is $A^\lambda$, so its value there is
 $-z-\epsilon\pi i$, whereas
 \[
  \FP_x\int_0^1 R_z(x)^{-1}\,dx
   =\FP_x\int_0^1(K(x)-z)\,dx=-z.
 \]
 The last equality follows by integrating $K=(\log\sin\pi x)'$
 on the symmetric cutoff interval. Thus the positive-resonance
 rule cannot simply be copied to negative powers.

 \item \textbf{Rational kernel exponents.}
 For a fixed nonintegral rational $\lambda$, determine whether
 $\mathcal L_\lambda(u,q)$ admits a finite representation in a
 specified hypergeometric or multiple-polylogarithm class at
 integral $u$. Its Gamma specialization at $q=0$ and its exact
 exponent jets are constraints on any proposed reduction. The
 unrestricted complex-$\lambda$ Dirichlet series is presently
 retained as an explicit coordinate.

 \item \textbf{The arithmetic of $\beta_1$.}
 Decide whether the normalized integral in
 \eqref{bhr:quartic-harmonic:beta-mellin} reduces to a specified
 algebra generated by ordinary zeta derivatives, Stieltjes
 constants, and Gamma or Barnes values at rational arguments.
 The proved identity for $Q_4+4\gamma Q_3$ gives an exact
 equivalence with a digamma finite-part evaluation, not an
 arithmetic independence statement. A new reduction must add
 information beyond that equivalence.

 \item \textbf{Higher symmetric Tornheim coefficients.}
 Use \eqref{bhr:fcy:Jintegral} to extract the three homogeneous cubic
 coefficients of $J_3$, including their elementary corrections.
 They control the cyclic sixth-derivative layer. Determine which
 reflection or differential identities reduce the formal
 three-coordinate count. For the present quadratic layer, seek
 a second exact representation of $\chi$ that could lead to a
 finite evaluation of $\chi$ or $\Omega_5$.

 \item \textbf{Commensurable Barnes periods.}
 Extend the individual resolvent theorem from unit periods to
 positive rationally commensurable periods. Multiplicity
 sequences should split into finitely many polynomial
 residue-class pieces, suggesting shifted Hurwitz and root-of-unity
 polylogarithm transforms. Prove that reduction with the exact
 normalization; arbitrary incommensurable periods are a separate
 problem and need not remain at finite depth two.

 \item \textbf{Rational shifts and cyclotomic colors.}
 Combine the polynomial weight with a rational Hurwitz shift,
 or with a finite Fourier character. Determine the exact
 colored analogue of $\mathcal O_{\alpha,\beta}(q)$ and its
 residue cancellation. An effective theorem should specify
 cumulative color products and convergence, since individual
 arguments of modulus greater than one need not obstruct the
 coupled nested sum.

 \item \textbf{Exact certificates for the current $S_6$ and $S_8$.}
 Keep each current candidate vector fixed. Try to express its
 difference from the claimed integral in a finite basis of
 iterated integrals or a rigorously defined cyclotomic quotient,
 and reduce that expression by proved shuffle, stuffle,
 distribution, and path relations. The deliverable should be
 an exact certificate, or an exact obstruction to the specified
 reduction, rather than additional matching digits. The
 already-rejected old $S_8$ vector is not the current target.

 \item \textbf{Formal finite-part and convergence infrastructure.}
 Formalize a reusable theorem that subtracts a parameter-dependent
 endpoint jet, integrates its monomials, and compares coefficient
 extraction with the original-coordinate finite part. Pair it
 with normal-convergence lemmas for a geometric outer index and
 a subtracted power-law inner index. These two modules would
 support most identities here without treating each moving pole
 as a separate ad hoc calculation.
\end{enumerate}

\subsection{Integration guide}

The suggested integration order is the polynomial transform,
the individual Barnes theorem, the complex-power generator,
the quartic harmonic bridge, and finally the transverse Tornheim
generator. The first two share their new analytic input;
the remaining branches have their own local subtraction proofs.
All manuscript labels in the deliverable use a common prefix.
The standalone source and the modular source compile the same
article. No external PDFs are bundled: the bibliography links to
the primary sources, while the provenance records identify the
exact repository snapshot and the reused numerical evaluator.

% END sections/07_audit_research.tex


% BEGIN references.tex
\begin{thebibliography}{99}

\bibitem{bhr:ProveIt}
V.~Reshetnikov, \emph{ProveIt: polylogarithm manuscript and incoming
research collection}, repository snapshot
\texttt{7bd45777a8a127e5dc69aff8887ca36a79a3059d},
11 October 2026.
\href{https://github.com/VladimirReshetnikov/ProveIt/tree/7bd45777a8a127e5dc69aff8887ca36a79a3059d/Analysis/Polylogarithms/docs/manuscript}{Pinned manuscript};
\href{https://github.com/VladimirReshetnikov/ProveIt/tree/7bd45777a8a127e5dc69aff8887ca36a79a3059d/docs/incoming}{pinned incoming collection}.
The attached inventory lists all source files used for the targeted audit.

\bibitem{bhr:Parity}
ProveIt research collection,
\emph{Harmonic Parity and Resolvent Identities},
incoming report dated 11 October 2026, in \cite{bhr:ProveIt}.
Sections 5 and 7; in particular the local remainder and symmetric
germ with labels \src{qtr:polar}, \src{qtr:kernel}, \src{qtr:J},
and questions R4, R6, R10.
\href{https://github.com/VladimirReshetnikov/ProveIt/blob/7bd45777a8a127e5dc69aff8887ca36a79a3059d/docs/incoming/ProveIt_Harmonic_Parity_and_Resolvent_Identities_2026-10-11.zip}{Pinned source archive}.

\bibitem{bhr:Cubic}
ProveIt research collection,
\emph{Cubic Tornheim Derivatives and Harmonic Stieltjes Constants},
incoming report dated 11 October 2026, in \cite{bhr:ProveIt}.
The cubic digamma--harmonic bridge and Section 7, question 6.
\href{https://github.com/VladimirReshetnikov/ProveIt/blob/7bd45777a8a127e5dc69aff8887ca36a79a3059d/docs/incoming/ProveIt_Cubic_Harmonic_Mixed_Orders_2026-10-11.zip}{Pinned source archive}.

\bibitem{bhr:DLMFZ}
NIST Digital Library of Mathematical Functions,
\S25.11, \emph{Hurwitz Zeta Function}.
\url{https://dlmf.nist.gov/25.11}.
Fourier expansion, argument shifts, Bernoulli values, and derivatives.
Accessed 11 October 2026.

\bibitem{bhr:DLMFG}
NIST Digital Library of Mathematical Functions,
\S\S5.7, 5.11, 5.15, 5.17, \emph{Gamma Function}.
\url{https://dlmf.nist.gov/5.15};
\url{https://dlmf.nist.gov/5.17}.
Digamma and polygamma expansions, asymptotics, and the Barnes $G$ function.
Accessed 11 October 2026.

\bibitem{bhr:DLMFPolylog}
NIST Digital Library of Mathematical Functions,
\S25.12, equation 25.12.12, \emph{Polylogarithms}.
\url{https://dlmf.nist.gov/25.12.E12}.
The Jonqui\`ere expansion used in Mellin subtractions.
Accessed 11 October 2026.

\bibitem{bhr:EMIntegrals}
O.~Espinosa and V.~H.~Moll,
\emph{On some integrals involving the Hurwitz zeta function: Part 2},
Ramanujan Journal \textbf{6} (2002), 449--468.
\url{https://www.math.tulane.edu/~vhm/papers_html/hurwitz2.pdf}.

\bibitem{bhr:EMGeneralized}
O.~Espinosa and V.~H.~Moll,
\emph{A generalized polygamma function},
Integral Transforms and Special Functions \textbf{15} (2004), 101--115.
\url{https://arxiv.org/abs/math/0305079}.

\bibitem{bhr:Adamchik}
V.~S.~Adamchik,
\emph{Multiple Gamma Function and Its Application to Computation of Series},
arXiv:math/0308074 (2003).
\url{https://arxiv.org/abs/math/0308074}.
Classical multiple-Gamma normalization and Hurwitz-zeta conversion.

\bibitem{bhr:Vardi}
I.~Vardi,
\emph{Determinants of Laplacians and multiple gamma functions},
SIAM Journal on Mathematical Analysis \textbf{19} (1988), 493--507.
\url{https://doi.org/10.1137/0519035}.

\bibitem{bhr:Mezo}
I.~Mez\H{o},
\emph{The Fourier series of the log-Barnes function},
arXiv:1604.00753v1 (2016).
\url{https://arxiv.org/pdf/1604.00753v1}.
Known Fourier coefficients; the pointwise correction in
Section~\ref{bhr:sec:audit} concerns equation (7), printed page 6 of this
specific version, and preserves its integrated use.

\bibitem{bhr:Kargin}
L.~Karg\i n, A.~Dil, M.~Cenkci, and M.~Can,
\emph{On the Stieltjes constants with respect to harmonic zeta functions},
Journal of Mathematical Analysis and Applications \textbf{525} (2023),
article 127302.
\url{https://arxiv.org/abs/2304.03517};
\url{https://doi.org/10.1016/j.jmaa.2023.127302}.
The harmonic-zeta coefficient framework is prior work; our strict
depth-three normalization is specified independently in the text.

\bibitem{bhr:Nakamura}
T.~Nakamura,
\emph{Rapidly convergent series representations of symmetric Tornheim
double zeta functions}, arXiv:2103.16873 (2021), Theorem 1.2.
\url{https://arxiv.org/abs/2103.16873}.

\end{thebibliography}

% END references.tex

\end{document}
